% 论文入口。沿用原 cumcmthesis.cls / cumcm2026.sty；原 example.tex 未覆盖。
% 在本目录运行 latexmk -xelatex main.tex 编译。
\documentclass[withoutpreface,bwprint]{cumcmthesis}
\usepackage{cumcm2026}
\usepackage{url}
\usepackage{placeins}
\hypersetup{hidelinks}
\setlength{\emergencystretch}{2em}
\renewcommand{\arraystretch}{1.13}
% 一级标题沿用中文数字；二、三级标题显示为4.1、4.2.1。
\renewcommand{\thesection}{\chinese{section}、}
\renewcommand{\thesubsection}{\arabic{section}.\arabic{subsection}}
\renewcommand{\thesubsubsection}{\thesubsection.\arabic{subsubsection}}
\title{非对称结算机制下的微网采购与储能协同调度}
\tihao{C}
\yearinput{2026}
\monthinput{9}
\dayinput{10}

\begin{document}
\maketitle
\begin{abstract}
\setlength{\parskip}{0.45em}
负载与光伏出力的波动使微网购电面临双重代价：采购不足需要高价补购，采购富余仍须按计划量付费；储能则受容量、功率和效率限制。为协调购电承诺与储能运行，本文以实际结算费用最小化为目标，建立统一调度框架，依次研究确定性供需、预测误差、日内信息更新和动态电价下的采购决策。

针对问题一，在负载与光伏已知的条件下，建立\textbf{\heiti 购电与储能联合线性规划模型}。将功率换算为时段电量，以供给不小于需求为约束，联合优化购电和充放电，并利用免费弃余条件处理充放电互斥。求得全天购电59482.70 kWh，购电费\textbf{35126.95元}，较无储能基准\textbf{节省26.90\%}，首末储电量均为6000 kWh。

针对问题二，面对日前预测误差及五倍价格的紧急补购，构建\textbf{\heiti 基于整日历史残差的场景采购模型}。先按一月实际费用校准预测器与储能执行器，再保留此前28日残差轨迹生成场景，优化共同购电计划，并根据实际供需逐段控制储能。在2月至12月334天回测中，总费用为\textbf{1401.28万元}，较相同预测及初态下的点预测采购\textbf{节省8.36\%}。

针对问题三，利用日内发布的新预报和实测储电量，建立\textbf{\heiti 双时间尺度滚动调度模型}，并引入\textbf{\heiti 基于执行规则回放的购电候选筛选}。在0、6、12、18时生成场景、点预测和误差余量采购候选，日内另保留当前合同；以相同情景、初态和执行规则回放评分，再按每10分钟的实际供需执行储能。四种筛选时点配置中，\textbf{\heiti 零时与日内均筛选}的费用最低，为\textbf{1387.15万元}，较未筛选母模型\textbf{降低0.747\%}。

针对问题四，考虑电价与净负荷的共同波动，构建\textbf{\heiti 同日配对的价格与净负荷联合场景模型}，保留高价格与高缺口同时出现的损失。日前版以场景LP制定计划，滚动版沿用全时点候选筛选，并按各情景对应价格计算回放费用、按实际价格结算。两版总费用分别为\textbf{1479.41万元}和\textbf{1467.02万元}，滚动版较其未筛选母模型\textbf{降低0.670\%}。

本文将非对称结算、储能物理约束与固定执行规则回放相结合，形成\textbf{\heiti 统一母模型与滚动采购筛选相结合的分问策略}。筛选适用问题与启用时点由开发期回测确定，日前两组保留母模型；滚动两组的费用下降伴随紧急补购增加，体现经济性与补购需求的权衡，不能据此保证未来费用下降。

\keywords{微网调度\quad 非对称结算\quad 场景优化\quad 执行回放\quad 候选筛选}
\end{abstract}

\section{问题重述}

光伏微网需要在负载需求、可再生出力和外部电网收费之间安排购电与储能。电池可以把低价时段购入或暂时富余的电量转移到后续时段，但其容量、功率和充放电效率限制了调节能力。题目要求在满足负载的基础上尽可能节省购电费用，因此应把真实结算账单作为比较策略的主要依据。

四个问题逐步改变可用信息和决策机会。问题一给出确定性输入，求单日日前购电及储能调度；问题二面对未知负载和光伏，需要先预测再制定日计划，缺口通过五倍价格紧急补购；问题三在规定预报时点允许付费修订常规购电，并根据实测储电量控制电池；问题四进一步引入价格不确定性，分别推广问题二的静态采购和问题三的滚动调度。

全文按问题顺序展开，每问依次讨论问题分析、模型建立和求解结果。以确定性LP和非对称结算下的场景采购构成统一母模型，再在滚动采购中加入基于执行规则回放的候选筛选；通过筛选时点消融选择最终配置，并报告实际费用与紧急补购的权衡。问题四进一步使用价格与净负荷的联合场景。每一问先解释新增机制，再给出变量、目标与约束，最后结合实际执行结果和对照实验判断模型效果。

\section{模型假设}

\begin{enumerate}
 \item 将附件中的功率解释为以标记时刻为右端点的10分钟平均功率，以$\Delta t=1/6$小时换算电量。这是时间离散的建模解释，不是题目明确给出的物理事实。
 \item 充电、放电效率分别为$\eta_c=\eta_d=0.9$，往返效率为0.81；忽略线路损耗和自放电。电池外侧充放电功率均不超过5000 kW，储电量保持在1200--10800 kWh。
 \item 不向外网售电，无法利用的富余电量可免费弃置。主结算口径中原计划全额付费且取消部分不退款；另设退款解释作敏感性对照。题目未给出的寿命折旧、碳成本或投资费用不进入主目标。
 \item 每10分钟可以观测当期负载、光伏、电价及储电量并完成平衡。紧急购电补足剩余缺口，不另设题目未给出的外网容量上限；未来实际值及尚未发布的预报不能提前用于决策。
 \item 问题一首末储电量均为6000 kWh。其他问题仅在1月1日初始化，之后实际状态逐日连续传递，不在每天零时或2月1日重新充满电池。
\end{enumerate}

\section{符号说明与数据准备}

\subsection{主要符号与统一单位}

购电、充放电和储电量均以kWh计，电价以元/kWh计。以$p$表示电价、$c$表示充电量，避免成本与充电变量混用。跨日问题增加日期下标$D$，不致混淆时省略。

\begin{table}[htbp]\centering\small
\caption{主要符号及含义}\label{tab:symbols}
\begin{tabularx}{\linewidth}{p{0.19\linewidth}Xl}\toprule
符号 & 含义 & 单位\\\midrule
$D,t,s$ & 日期、时段$t=0,\ldots,143$、场景索引 & ---\\
$L_{D,t},R_{D,t}$ & 实际负载与光伏功率 & kW\\
$n_{D,t},p_{D,t}$ & 实际净负荷电量、实际交易电价 & kWh、元/kWh\\
$g_{D,t},a_{D,t}$ & 零时计划、最终生效常规购电量 & kWh\\
$c_t,d_t,e_t$ & 电池外侧充电量、放电量、时段初储电量 & kWh\\
$h_t,w_t$ & 紧急购电量、弃置电量 & kWh\\
$\pi_s,Q_s$ & 场景概率、场景剩余时域紧急费用 & ---、元\\
$\rho_t,\lambda$ & 参考备用量、CVaR权重 & kWh、---\\
$\beta_T,k_\rho$ & 终端偏离惩罚系数、备用缩放系数 & 元/kWh、---\\
$E_{\min},E_{\max},U$ & 储电量下上限、单段充放电量上限$5000/6$ & kWh\\
\bottomrule\end{tabularx}
\end{table}

\subsection{统一物理时间与模板列序}

每日保留144个10分钟时段，物理区间及净负荷定义为
\begin{equation}
 I_t=[t/6,(t+1)/6),\qquad n_{D,t}=(L_{D,t}-R_{D,t})\Delta t.
 \label{eq:time_energy}
\end{equation}
源标记00:10对应物理00:00--00:10，最后一个源标记次日00:00对应当日23:50--24:00。模型按物理自然日结算，而附件5购电页保留原始首标签0:10--0:20及末标签次日0:00--0:10，按源样本列序写入，不根据显示标签再次平移结果。

\begin{table}[htbp]\centering\small
\caption{计算时间与原模板显示的对应}\label{tab:physical_mapping}
\begin{tabularx}{\linewidth}{lLLL}\toprule
索引 & 源采样右端点 & 模型物理区间 & 原模板显示区间\\\midrule
$t=0$ & 00:10 & 00:00--00:10 & 0:10--0:20\\
$t=60$ & 10:10 & 10:00--10:10 & 10:10--10:20\\
$t=143$ & 次日00:00 & 23:50--24:00 & 次日0:00--0:10\\
\bottomrule\end{tabularx}
\end{table}

144段映射冻结于\texttt{data/time\_mapping.json/csv}，同时记录各工作簿原标签及单元格位置。题目指定10:00--10:10取$t=60$；四小时区间累加连续24段；同日相邻正紧急量合并为一个区间。这些操作均调用同一物理索引，热力矩阵与能量图也沿用该索引。

\subsection{评价数据及图表口径}

问题一使用附件1。问题二至四的评价和填报覆盖2025年2月1日至12月31日334天，一月用于历史积累与有限候选校准。各问题的具体预测及可用信息分别在对应模型中说明。同一问题的比较策略从共同2月初态出发，各自连续执行；不同问题保留各自已冻结一月运行形成的初态。

预测器与实时执行器沿用原冻结配置；候选筛选的适用问题及启用时点依据开发性回测确定。最终配置确定后固定，本文所报告结果属于开发期闭环评价。第二问与第四问日前版保留未筛选母模型，第三问与第四问滚动版在零时及全部合法日内修订时点启用筛选。决策过程只使用当时可得信息，与选型参考过开发期费用是两个不同层面的要求，不能据前者宣称未接触数据上的独立测试。

实际总费用只计算计划、调整及紧急三项账单，辅助终端罚项和风险权重不加入报表。正文表格和附件5均从同一执行数组生成，保留全部指定日期和逐日数据。新增热力图展示时序分布或策略差异，瀑布图分解实际费用变化；分位阴影带表示场景或历史窗口间差异，不作为未经估计的置信区间。

\FloatBarrier
\section{问题一模型建立与求解}

\subsection{问题一分析}

问题一给定单日负载、光伏预测及电价，并要求在零时一次性形成144个10分钟时段的购电计划，日末储电量恢复到日初水平。其关键是决定电量在同一时段内如何分配，以及通过电池向后续时段转移多少电量。低价购电未必立即用于负载，光伏富余也可补充储能；相反，高价时段是否放电取决于当时缺口、剩余电量及后续补能机会。

本文把附件1的已给曲线视为本问确定输入，先求解购电与储能的联合线性规划，为后续不确定条件下的采购提供共同物理基础。已有日前微网调度研究\cite{kong2016,silva2020}支持这一建模结构，本题只保留外网、光伏、负载和储能，不引入题目未给出的其他发电设备、售电收入或寿命折旧。确定性结果解释的是在给定输入下如何调度，不能作为未知实际负载时已经掌握完美信息的证明。

\subsection{问题一模型建立}

\subsubsection{决策变量与购电目标}

按统一物理索引$t=0,\ldots,143$，定义计划购电$g_t$、电池外侧充电$c_t$、电池外侧放电$d_t$和时段初储电量$e_t$，单位均为kWh。净负荷为$n_t=(L_t-R_t)\Delta t$，可以为负，负值表示光伏超过当期负载。购电按时段电价$p_t$收费，因此目标为
\begin{equation}
 \min_{g,c,d,e}\quad C_1=\sum_{t=0}^{143}p_tg_t.
 \label{eq:q1_objective}
\end{equation}
这里$g_t$已是10分钟电量，费用计算不再乘一次$\Delta t$。目标只含赛题要求支付的外购电费；免费弃余、已有电池和光伏均不额外产生本问未给定的收费项目。

\subsubsection{供需、储能与日末约束}

在任一时段，外网采购与电池放电至少补足净负荷及充电需求：
\begin{equation}
 g_t+d_t-c_t\ge n_t.
 \label{eq:balance}
\end{equation}
令$w_t=g_t+d_t-c_t-n_t\ge0$，即可改写为电量平衡等式。$w_t$表示未消纳的富余量，不能解释为向外网售电。

电池状态沿时间传递：
\begin{equation}
 e_{t+1}=e_t+\eta_cc_t-\frac{d_t}{\eta_d},
 \qquad \eta_c=\eta_d=0.9.
 \label{eq:soc}
\end{equation}
充电量$c_t$在电池外侧计量，只有$0.9c_t$进入电池；要对外输出$d_t$，内部需要消耗$d_t/0.9$。因此不能将充、放电项都乘0.9，也不能把完整往返效率误当为90\%。

容量和功率约束为
\begin{equation}
 E_{\min}\le e_{t+1}\le E_{\max},\qquad
 0\le c_t,d_t\le U,
 \quad E_{\min}=1200,\quad E_{\max}=10800,\quad U=5000/6.
 \label{eq:limits}
\end{equation}
储电量下限防止过度放电，上限防止过度充电；$U$由5000 kW乘$1/6$小时得到，分别限制一个时段可输入、输出的外侧电量。另有
\begin{equation}
 g_t\ge0,\qquad e_0=e_{144}=6000.
 \label{eq:terminal_q1}
\end{equation}
非负购电排除反向售电，首末相等排除通过额外消耗初始库存制造的单日账面节费。模型由线性目标、线性等式和不等式组成，可直接作为LP求解。

\subsubsection{免费弃余条件下的充放电互斥证明}

通常可用二进制变量禁止同时充放电，但本题免费弃余条件使其并非必要。若某可行解在时段$t$有$c_t,d_t>0$，取
\begin{equation}
 \varepsilon=\min\{c_t,d_t/(\eta_c\eta_d)\},\qquad
 c'_t=c_t-\varepsilon,\quad
 d'_t=d_t-\eta_c\eta_d\varepsilon.
 \label{eq:q1_remove_simultaneous}
\end{equation}
变换后至少有一种动作变为零，充放电量均不增大，因此不破坏功率上下界。同时
\[
 \eta_cc'_t-d'_t/\eta_d=\eta_cc_t-d_t/\eta_d,
\]
故全部储电轨迹与首末约束保持不变。供需约束左端增加$(1-\eta_c\eta_d)\varepsilon\ge0$，可将新增富余弃置，原购电量及费用无须增加。对所有同时充放电时段依次变换，得到满足互斥且费用不高于原解的方案。因此至少存在一组充放电互斥的最优解，无需为该条件额外引入整数变量。

程序在充放电总量上加入$10^{-7}$元/kWh的数值破同解项，以减少求解器返回无意义循环的可能；统计实际费用时剔除该项。以上证明依赖允许免费弃余，不能不加条件地推广到存在弃电罚金或必须消纳的其他问题。

\subsubsection{跨时段储能的经济条件}

若低价时段从外网购电并储存，要在高价时段对外输出1 kWh，低价段需购入$1/(\eta_c\eta_d)$ kWh。因此在忽略容量、功率及其他机会成本时，纯购电套利的必要盈利条件为
\begin{equation}
 p_h>\frac{p_\ell}{\eta_c\eta_d}=\frac{p_\ell}{0.81}.
 \label{eq:arbitrage}
\end{equation}
该式解释电池为何偏向低价补能、高价替代购电，但不足以单独确定完整计划。例如预期午间光伏可充满电池时，凌晨继续充电可能占用本可容纳免费光伏的容量；临近日末则还须为6000 kWh终值补能。完整LP同时处理这些跨时段竞争。

\subsection{问题一求解}

\subsubsection{求解步骤与结果核对}

首先按照统一时间映射读取附件1，将功率乘$\Delta t$得到净需求及单段充放电上限；然后以式\eqref{eq:soc}组装相邻状态的稀疏约束矩阵，加入购电、功率、容量和首末约束，由HiGHS求解。最后从未四舍五入的变量计算费用与汇总电量，核对144段供需平衡、储电递推、容量和功率界以及单段互斥。工作簿按原模板的源样本列序填写，论文指定时段则通过同一个物理索引取数。

计算得到全天计划购电量59482.6990 kWh、购电费35126.9486元。无储能基准逐时购买$[n_t]_+$，费用为48052.05元，因此采用储能节省12925.10元，节费率为26.90\%。电池外侧累计充电20740.67 kWh、放电16799.94 kWh，满足
\[
 \eta_c\sum_tc_t=\sum_td_t/\eta_d,
\]
与$e_0=e_{144}=6000$ kWh一致。充放电外侧累计量不相等，差额来自效率损耗，不是储能平衡错误。

\begin{table}[htbp]\centering\small\caption{问题1指定时段与全天购电结果（电量单位：kWh）}\label{tab:q1-purchase}\begin{tabularx}{\linewidth}{lRlRlR}
\toprule
时间段 & 购电量 & 时间段 & 购电量 & 时间段 & 购电量\\
\midrule
10:00--10:10 & 0.00 & 12:00--12:10 & 480.41 & 14:00--14:10 & 0.00\\
16:00--16:10 & 445.43 & 18:00--18:10 & 531.89 & 20:00--20:10 & 0.00\\
\midrule
全天购电量 & 59482.70 & \multicolumn{2}{c}{全天购电费（元）} & \multicolumn{2}{r}{35126.95}\\
\bottomrule
\end{tabularx}\end{table}
\begin{table}[htbp]\centering\small\caption{问题1储能充放电及首末电量（单位：kWh）}\label{tab:q1-battery}\begin{tabularx}{\linewidth}{lRRlRR}
\toprule
时间段 & 充电量 & 放电量 & 时间段 & 充电量 & 放电量\\
\midrule
00:00--04:00 & 4500.00 & 0.00 & 04:00--08:00 & 833.33 & 6365.84\\
08:00--12:00 & 4787.96 & 1703.00 & 12:00--16:00 & 5286.04 & 91.10\\
16:00--20:00 & 0.00 & 5780.13 & 20:00--24:00 & 5333.33 & 2859.87\\
\midrule
0:00储电量 & \multicolumn{2}{r}{6000.00} & 24:00储电量 & \multicolumn{2}{r}{6000.00}\\
\bottomrule
\end{tabularx}\end{table}

表\ref{tab:q1-purchase}中的全天购电量包含为负载和储能安排的所有外购电，表\ref{tab:q1-battery}则按六个物理四小时区间汇总充放电。某个四小时区间内两项均为正，表示其中不同10分钟时段分别充电或放电，不违反单时段互斥。

\subsubsection{调度结构与能量收支}

图\ref{fig:q1}将购电、充放电、SOC及电价对齐展示，用来判断电池如何在价格和净需求变化时转移电量。原始LP未包含爬坡、启停或平滑约束，因此保留其分段动作，不对最优结果额外平滑后再声称仍然最优。

\begin{figure}[htbp]\centering
\includegraphics[
    width=0.90\linewidth,
    height=0.60\textheight,
    keepaspectratio
    ]{figures/results/q1_dispatch.pdf}
\caption{问题一购电、充放电、储电量及电价的日内对应关系}
\label{fig:q1}
\end{figure}

图\ref{fig:q1-energy-heatmap}把六个四小时区间的购电、光伏、负载、充电和放电量对齐展示，格内标注原始累计电量，便于与题目汇总表核对。下方平均库存另用独立色标，按区间内24段$(e_t+e_{t+1})/2$的平均计算，不与流量混作同一统计量。

\begin{figure}[htbp]\centering
\includegraphics[width=0.94\linewidth]{figures/results/q1_energy_heatmap.pdf}
\caption{问题一六个物理四小时区间的累计电量与平均库存；流量共用色标，库存另设色标，标注单位均为kWh}
\label{fig:q1-energy-heatmap}
\end{figure}
\FloatBarrier

\FloatBarrier
\section{问题二模型建立与求解}

\subsection{问题二分析}

问题二在每日零时承诺整天购电量，但真实负载和光伏尚未知晓。计划量偏少时须支付五倍交易电价补足，偏多时已付费用不能收回。因而把一条均值预测直接代入问题一，未必能得到期望费用最低的采购量。储能同时参与价格套利和误差补偿，某时段的预测误差还会改变后续可用电量，必须按完整轨迹评估。

本问采用“因果预测、整日残差场景、共同采购计划、实际状态反馈”的方案\cite{abunima2022,marino2018,grover2018}。母模型预先关闭额外CVaR及备用，即$\lambda=0$、$\rho_t=0$，场景数保留此前最多28条历史轨迹。预测器和执行器在一月验证中按实际账单选择，风险偏好与场景压缩只做后续对照。本问的候选筛选在开发期未取得节费收益，因此最终仍采用母模型；滚动问题的筛选选型另见第三问。

\subsection{问题二模型建立}

\subsubsection{五倍补购费用为何导出80\%采购分位数}

先去掉储能和跨时段耦合，研究单时段基准。设$X$为实际净需求电量，$q\ge0$为计划采购，固定电价$p>0$，富余无残值，且$\mathbb E[X_+]<\infty$。概率和期望均以决策时已知信息为条件。期望费用为
\begin{equation}
 J_2(q)=pq+5p\,\mathbb E[(X-q)_+].
 \label{eq:mech_q2_cost}
\end{equation}
第一项是无论是否使用都要支付的计划费，第二项是计划不足时的补购费。记$F_X(x)=\Pr(X\le x)$、$F_X(x^-)=\Pr(X<x)$。$J_2$为凸函数，左右导数给出的次梯度为
\begin{equation}
 \partial J_2(q)=p[5F_X(q^-)-4,\ 5F_X(q)-4].
 \label{eq:mech_q2_subgradient}
\end{equation}
对于$q>0$的内点，最优解的充要条件是零属于其次梯度，即
\begin{equation}
 F_X(q^-)\le0.8\le F_X(q).
 \label{eq:mech_q2_quantile_condition}
\end{equation}
若分布在最优点连续，则$J_2'(q)=p-5p\Pr(X>q)=0$，化为$F_X(q)=0.8$；离散场景存在概率跳跃时，使用式\eqref{eq:mech_q2_quantile_condition}，不强求累积概率恰好等于0.8。若最优点位于$q=0$边界，则右导数非负要求$F_X(0)\ge0.8$。以下分位数$F_X^{-1}(\alpha)=\inf\{x:F_X(x)\ge\alpha\}$表示，可选取
\begin{equation}
 q_2^\star=\max\{0,F_X^{-1}(0.8)\}.
 \label{eq:mech_q2_quantile}
\end{equation}
平坦分布区间可能导致多个最优解，上式给出其中一个。

这说明采购高于均值可能只是最小化期望费用的结果，无须额外诉诸风险厌恶。80\%源于计划价格与五倍补购价格的比例；加入储能后还存在效率、功率、容量及跨时段机会价值，该阈值只作机制解释和简化基准，不能直接逐时套用，也不是后文备用误差分位数的理论依据。

\subsubsection{因果预测与决策导向选型}

问题二只使用附件1电价和附件2历史负载、光伏，不使用附件3发布预报。候选负载预测为前一周、最近2个或4个同星期日的同一时段均值，光伏预测为此前3、7或14个完整日的同一时段均值。以$X$表示负载，同星期预测写为
\begin{equation}
 \widehat X_{D,t}=\frac{1}{|\mathcal W_D|}\sum_{j\in\mathcal W_D}X_{j,t},\qquad
 \mathcal W_D=\{D-7i:i=1,\ldots,K\}\cap\{1,\ldots,D-1\}.
 \label{eq:weekday_forecast}
\end{equation}
光伏移动均值改用此前$K$个完整日。历史不足时利用已有样本，没有相应星期日则退回近7日均值，首日采用附件1固定先验。为其他问题生成的历史价格预测可使用同一形式；每个历史预测都只使用其日期以前的数据，不能把实际未来曲线作为预测。

文献\cite{elmachtoub2017spo}提出按预测引致的下游决策损失评价预测。本题借鉴这一思想，在固定相同场景和控制器的条件下比较候选预测器的实际结算费用，而不只按RMSE选择。未采用SPO或SPO+训练算法，也不引述其训练理论为本程序保证。

验证窗口为1月15--31日。同一问题的所有候选均从冻结旧版策略在1月15日的实际电量7081.8620 kWh开始，逐日因果运行，选择分数为计划费加紧急费，不含终端罚项。先固定28场景、无风险及备用的内层LP，比较负载方法，再固定选中的负载方法比较光伏；随后固定预测器比较贪心与内层LP。该分步搜索并未穷举预测和控制的全部组合。

% Generated by code/make_question_tables.py; do not edit values.
% Source: data/selection_v2.json; SHA256=f290084567f8b335dfbe7c042da7bf6863af20d3db001376f8bb5003dedd9aa7
\begin{table}[htbp]\centering\small
\caption{问题二预测器的一月实际费用验证}
\label{tab:q2-forecast-validation}
\begin{tabularx}{\linewidth}{LLLR}\toprule
比较步骤 & 负载预测 & 光伏预测 & 实际费用/元\\\midrule
负载比较 & 前一周 & 最近7日 & 933430.35\\
负载比较 & 同星期2日 & 最近7日 & 929216.46\\
负载比较 & 同星期4日 & 最近7日 & 930211.02\\
光伏比较 & 同星期2日 & 最近3日 & 928995.30\\
光伏比较 & 同星期2日 & 最近7日 & 929216.46\\
光伏比较 & 同星期2日 & 最近14日 & 932521.33\\
\bottomrule\end{tabularx}
\par\smallskip\parbox{\linewidth}{\footnotesize 此处固定内层LP及其余设置；相同配置复用同一次计算。随后单独比较执行器。费用不含终端目标项。}
\end{table}


由表\ref{tab:q2-forecast-validation}确定负载采用最近2个同星期日均值、光伏采用近3日均值。在这一预测组合下，贪心执行的一月费用917318.59元，内层LP为928995.30元，故主案选择贪心规则。两者一月验证末库存分别7155.31和6000.00 kWh，未强制相等，不能把费用比较表述为库存价值已经中和。

2月初态共同承接已冻结旧版1月实际末态8875.1909 kWh，不用一月底选出的新方法事后回写1月。验证轨迹用于选型，旧版一月运行轨迹用于提供共同前缀，两者不能混淆。2月以后预测器类型、实时控制器及母模型设置固定，实际SOC连续传递；候选筛选的适用问题与启用时点依据2--12月开发性回测确定。因此，本文报告开发期闭环评价，不称为完全未见数据上的独立测试。

\subsubsection{整日历史残差与场景构造}

在日期$D$的零时，取此前至多28个完整日组成$\mathcal H_D$。历史日$j$在时段$t$的净负荷残差为
\begin{equation}
 \epsilon^n_{j,t}=
 \big[(L_{j,t}-R_{j,t})-(\widehat L_{j,t}-\widehat R_{j,t})\big]\Delta t,
 \qquad j<D.
 \label{eq:residual}
\end{equation}
每个历史日保留为完整144段轨迹，当前场景由
\begin{equation}
 n_{t,s}=\widehat n_t+\epsilon^n_{j_s,t},\qquad
 \pi_s=1/|\mathcal H_D|,\qquad \sum_s\pi_s=1
 \label{eq:q2_empirical_scenarios}
\end{equation}
给出。这样保留连续几小时共同偏高或偏低的误差，避免独立逐点抽样破坏时间关联。历史充足后排除采用固定先验的1月1日，主填报期均保留28条历史日轨迹。经验分布只代表有限窗口内出现过的误差，不等于真实未来分布。

场景缩减对照取最多12个代表日。对历史误差按整个残差块的总体标准差归一，计算轨迹间欧氏距离；先选总距离最小的代表，再逐次加入离已有代表最远的轨迹，随后交替分配样本与更新簇内medoid，至代表不变或达到8轮。每个代表仍是实际历史日，概率为所属簇频率。压缩改变场景及权重，其经济效果另行验证，不视为不影响结果的纯计算简化。

\subsubsection{受限两阶段采购模型}

日前第一阶段选择所有场景共享的计划$g_t$及参考$c_t,d_t,e_{t+1}$；补救阶段只允许紧急购电$h_{t,s}$随场景变化：
\begin{equation}
 g_t+d_t+h_{t,s}-c_t\ge n_{t,s},\qquad h_{t,s}\ge0.
 \label{eq:scenario_balance}
\end{equation}
同一组参考充放电须在所有场景下配合相应紧急补购满足净需求。以
$Q_s=\sum_t5p_th_{t,s}$表示一条完整场景的紧急费用，主目标为
\begin{equation}
 \min\left\{\sum_tp_tg_t+\sum_s\pi_sQ_s+
 \beta_T|e_{144}-6000|\right\},\qquad
 \beta_T=b_T\operatorname{median}_{t}p_t,\quad b_T=1.
 \label{eq:q2_objective}
\end{equation}
第一项是零时承诺的计划费，第二项按场景概率计算期望紧急费，二者体现本题的非对称结算。第三项是有限时域库存价值的软近似，抑制单纯为缩短本日账单而在日末过度放电的倾向，不计入实际账单，也不强制每天恢复6000 kWh。绝对值通过正负偏差变量线性化。容量、功率和状态方程沿用式\eqref{eq:soc}、\eqref{eq:limits}，初态使用当日实际SOC，购电非负。

该模型是受限两阶段随机规划，而非完全多阶段策略。共享储能参考避免对每个场景独立优化整天、提前知道未来情况；但实际控制会利用逐时观测调整储能，参考紧急费仍可能不同于真实执行费用。本文分别保留参考和实际轨迹，不把外层$c_t,d_t$直接作为附件5的实际充放电量。

\subsubsection{贪心执行器及状态连续性}

零时购电计划一经确定，本问全天令$a_t=g_t$。当期可观测的有效采购$a_t$、净需求$n_t$和储电量$e_t$决定实际动作。若$a_t\ge n_t$：
\begin{equation}
 c_t^{\rm act}=\min\{a_t-n_t,U,(E_{\max}-e_t)/\eta_c\},\quad
 d_t^{\rm act}=h_t=0,\quad w_t=a_t-n_t-c_t^{\rm act}.
 \label{eq:execute_charge}
\end{equation}
若$a_t<n_t$：
\begin{equation}
 d_t^{\rm act}=\min\{n_t-a_t,U,\eta_d(e_t-E_{\min})\},\quad
 c_t^{\rm act}=w_t=0,\quad h_t=n_t-a_t-d_t^{\rm act}.
 \label{eq:execute_discharge}
\end{equation}
然后按式\eqref{eq:soc}更新状态。此规则优先以常规富余充电、以当前电量补缺，不足部分紧急购买，不用紧急电量主动充电。充放电互斥、容量和功率约束逐步得到满足，外网补购使实际负载得到保障。

日际满足$e_{D+1,0}=e_{D,144}$，不每日重置。执行器不直接读取当前或未来价格、参考SOC和参考备用，因此其与场景规划的差异仍然存在。后续问题会用固定合同回放和内层LP专门检查这种差异；贪心被一月账单选中，不代表已经证明它是所有因果控制策略中的最优解。

\subsubsection{风险与备用的可选扩展}

经济主案取$\lambda=\rho_t=0$。为了识别额外保守程度的影响，对照在式\eqref{eq:q2_objective}加入$\lambda\operatorname{CVaR}_{0.90}(Q)$。其中$Q$只含紧急费用，不是全天总账单；按照文献\cite{rockafellar2000}写为
\begin{equation}
 \operatorname{CVaR}_{\alpha}(Q)=
 \min_z\left\{z+\frac{1}{1-\alpha}\sum_s\pi_sv_s\right\},\quad
 v_s\ge Q_s-z,\quad v_s\ge0,\quad\alpha=0.90.
 \label{eq:cvar}
\end{equation}
自适应权重根据近期残差总体标准差设定：
\begin{equation}
 \lambda=0.10\min\{3,r/0.20\},\qquad
 r=\frac{\operatorname{mean}_t\operatorname{sd}_j(\epsilon^n_{j,t})}
 {\max\{\operatorname{mean}_t|\widehat n_t|,50\}}.
 \label{eq:risk_weight}
\end{equation}
标准差分母取历史样本日数，50 kWh防止净需求接近零时比值失稳。固定风险权重则取一月自适应序列均值0.10368，不能用2--12月数据重新匹配。

参考备用沿同一历史日累加未来至多六段误差，以外侧累计缺口换算为内部库存：
\begin{align}
 \rho_t^{\rm raw}
 &=k_\rho\min\left\{2400,
 \frac{\left[q_{0.8}\left(\sum_{k=t+1}^{\min(t+6,143)}
             \epsilon^n_{j,k}\right)\right]_+}{\eta_d}\right\},
 \label{eq:reserve}\\
 \rho_t&=\min\{\rho_t^{\rm raw},
 [e_\tau-E_{\min}+\eta_cU(t-\tau+1)]_+\},\qquad
 e_{t+1}\ge E_{\min}+\rho_t.
 \label{eq:reachable_reserve}
\end{align}
本问$\tau=0$，后续滚动问题可取其他发布时点。分位数按排序样本间线性插值，空和为零；先在2400 kWh截断，再乘无量纲$k_\rho$，故缩放后的上限为$2400k_\rho$。最后按当前电量及最大充电速度做可达截断。该处理与先缩放再重新截到2400不同，复现时须保持顺序。

动态备用对照取$k_\rho=1$；固定备用取一月可达截断前动态备用均值208.67 kWh，日末置零，再作相同可达截断。一月匹配按每次外层更新实际覆盖的时段计数。备用仅限制外层参考轨迹，实际执行器可为当期负载消耗它，物理下限仍为1200 kWh，不能解释成保证80\%供电可靠性的硬约束。终端系数$\beta_T$单位为元/kWh，备用系数$k_\rho$无量纲，二者分开记号。

\subsection{问题二求解}

\subsubsection{逐日求解与评价口径}

每天零时先读取之前已完成的负载及光伏，更新冻结类型的预测器，再按式\eqref{eq:residual}构造28条场景；以当前实际SOC为初态求解外层LP，保存整日原计划和参考储能。随后每10分钟只揭示当前实际负载、光伏，按式\eqref{eq:execute_charge}或\eqref{eq:execute_discharge}执行，将本段充放电、紧急量及末态写入实际数组。当天结束后按
\begin{equation}
 C_D=\sum_{t=0}^{143}p_tg_{D,t}+\sum_{t=0}^{143}5p_th_{D,t}
 \label{eq:q2_actual_bill}
\end{equation}
重新结算，并把末态传至下一日。优化器的CVaR、终端项及数值破同解项都不进入该账单。

正式评价范围为2025年2月1日至12月31日334天、每天144段。除核对费用可加性外，还检查实际电量平衡、SOC递推、充放电互斥、上下界、日际连续性以及全天$a_t=g_t$。指定日表、日汇总和附件5均从同一执行数组生成，原模板表头和列序保持不变。

预测的RMSE及MAE只描述所选方法，不参与回测期再选型。表\ref{tab:q2-prediction-errors}给出实际误差，说明低成本采购仍建立在有误差的预测之上。
% Generated by code/make_question_tables.py; do not edit values.
% Source: data/selection_v2.json; SHA256=f290084567f8b335dfbe7c042da7bf6863af20d3db001376f8bb5003dedd9aa7
% Source: data/prepared_v2_q2.npz; SHA256=0322336272736ef1956aefa31f183bf42740abd1797e8d0d290fc61a700158fa
\begin{table}[htbp]\centering\small
\caption{问题二已选预测器的误差，单位为kW}
\label{tab:q2-prediction-errors}
\begin{tabularx}{\linewidth}{LRRRR}\toprule
变量 & 一月MAE & 一月RMSE & 评价期MAE & 评价期RMSE\\\midrule
负载 & 122.51 & 152.09 & 171.91 & 240.54\\
光伏 & 100.38 & 214.94 & 154.46 & 313.48\\
\bottomrule\end{tabularx}
\par\smallskip\parbox{\linewidth}{\footnotesize 一月指标取1月15--31日，评价期取2--12月。类型按一月结算费用选择，表中误差不用于重选。}
\end{table}


\subsubsection{总体结果与基础方案比较}

计划费用13507968.43元、紧急费用504839.24元，总费用14012807.66元；各分项分别四舍五入时可能与总数相差0.01元，汇总统一从原精度计算。全天计划购电22047537.09 kWh，紧急购电85017.69 kWh，后者涉及113天、617个10分钟时段。年末库存6240.32 kWh，低于共同期初8875.19 kWh，因此比较账单时须同时看库存，不称其为库存中性节费。

\begin{table}[htbp]\centering\small
\caption{问题2在相同预测与初态下的基础对照}
\label{tab:v2_baselines}
\begin{tabularx}{\linewidth}{LRRRR}\toprule
策略 & 总费/万元 & 紧急量/kWh & 日费CVaR95/元 & 末库存/kWh\\\midrule
无储能场景采购 & 1854.80 & 242340.86 & 80425.37 & 8875.19\\
点预测采购 & 1529.08 & 505585.75 & 100557.90 & 5845.40\\
场景期望费用（主案） & 1401.28 & 85017.69 & 70731.74 & 6240.32\\
\bottomrule\end{tabularx}
\par\smallskip\parbox{\linewidth}{\footnotesize 本表基于未启用候选筛选的母模型，各组保留原实验配置及数值，覆盖相同334天并连续推进状态；不作为最终筛选算法的同条件消融。}
\end{table}


表\ref{tab:v2_baselines}的三组使用相同预测和期初库存。点预测采购仅把预测曲线输入调度器，无储能场景采购则禁用充放电，其他经济设置保持一致。点预测费用15290840.75元，主案节省1278033.08元，即8.36\%；相较无储能场景采购18548033.07元，主案费用降低24.45\%。主要证据支持历史场景对非对称补购代价的处理及储能跨时段转移，而不是必须加入CVaR。紧急购电已支付高额代价但补足了负载，其发生不表示实际已经失供电。

图\ref{fig:q2-calendar-heatmap}按日期和物理时段联合展示紧急购电与时段末SOC，可定位缺口在一天中的位置，以及是否伴随库存接近下限。零紧急量留浅灰，正值采用幂次色阶突出稀疏小额事件，色条仍标原始kWh；库存使用独立的固定上下界色标。不能仅凭两个面板的颜色深浅把补购量和库存直接作等量比较。
\begin{figure}[htbp]\centering
\includegraphics[width=0.97\linewidth]{figures/results/q2_calendar_heatmap.pdf}
\caption{问题二334个日期与144个物理时段的紧急量和时段末SOC；紧急量正值采用$\gamma=0.4$幂次色标，颜色数值仍以kWh标示}
\label{fig:q2-calendar-heatmap}
\end{figure}

基于执行规则回放的候选筛选在本问得到费用14035379.65元，比母模型增加22571.98元，即0.161\%，紧急量也增至138015.18 kWh。因此本问不启用该模块。它在滚动采购中的适用性及筛选时点比较见第三问，不能由滚动组节费推断日前采购同样受益。

\subsubsection{题目指定日期与完整策略填报}

\begin{table}[htbp]\centering\small
\caption{问题2指定日期购电结果（沿用题目表1布局，电量单位：kWh）}
\label{tab:q2-purchase}
\par\smallskip\textbf{2025-03-20}\par\smallskip
\begin{tabularx}{\linewidth}{lRlRlR}
\toprule
时间段 & 购电量 & 时间段 & 购电量 & 时间段 & 购电量\\
\midrule
10:00--10:10 & 0.00 & 12:00--12:10 & 564.20 & 14:00--14:10 & 0.00\\
16:00--16:10 & 490.32 & 18:00--18:10 & 704.86 & 20:00--20:10 & 0.00\\
\midrule
全天购电量 & 67796.64 & \multicolumn{2}{c}{全天购电费（元）} & \multicolumn{2}{r}{41219.41}\\
\bottomrule
\end{tabularx}
\par\smallskip\textbf{2025-06-21}\par\smallskip
\begin{tabularx}{\linewidth}{lRlRlR}
\toprule
时间段 & 购电量 & 时间段 & 购电量 & 时间段 & 购电量\\
\midrule
10:00--10:10 & 0.00 & 12:00--12:10 & 0.00 & 14:00--14:10 & 0.00\\
16:00--16:10 & 296.43 & 18:00--18:10 & 275.52 & 20:00--20:10 & 0.00\\
\midrule
全天购电量 & 40036.30 & \multicolumn{2}{c}{全天购电费（元）} & \multicolumn{2}{r}{23986.38}\\
\bottomrule
\end{tabularx}
\par\smallskip\textbf{2025-09-23}\par\smallskip
\begin{tabularx}{\linewidth}{lRlRlR}
\toprule
时间段 & 购电量 & 时间段 & 购电量 & 时间段 & 购电量\\
\midrule
10:00--10:10 & 0.00 & 12:00--12:10 & 475.29 & 14:00--14:10 & 0.00\\
16:00--16:10 & 537.29 & 18:00--18:10 & 789.93 & 20:00--20:10 & 0.00\\
\midrule
全天购电量 & 69576.50 & \multicolumn{2}{c}{全天购电费（元）} & \multicolumn{2}{r}{46384.26}\\
\bottomrule
\end{tabularx}
\par\smallskip\textbf{2025-12-21}\par\smallskip
\begin{tabularx}{\linewidth}{lRlRlR}
\toprule
时间段 & 购电量 & 时间段 & 购电量 & 时间段 & 购电量\\
\midrule
10:00--10:10 & 0.00 & 12:00--12:10 & 970.92 & 14:00--14:10 & 0.00\\
16:00--16:10 & 851.12 & 18:00--18:10 & 744.23 & 20:00--20:10 & 0.00\\
\midrule
全天购电量 & 101008.53 & \multicolumn{2}{c}{全天购电费（元）} & \multicolumn{2}{r}{65454.44}\\
\bottomrule
\end{tabularx}
\par\smallskip\parbox{\linewidth}{\footnotesize 注：问题2与4-2列示计划购电量；问题3与4-3列示最终生效的调整后常规购电量。全天购电量不含紧急购电，全天购电费含计划、调整与紧急费用。}
\end{table}

\begin{table}[htbp]\centering\small
\caption{问题2指定日期储能结果（沿用题目表2布局，单位：kWh）}
\label{tab:q2-battery}
\par\smallskip\textbf{2025-03-20}\par\smallskip
\begin{tabularx}{\linewidth}{lRRlRR}
\toprule
时间段 & 充电量 & 放电量 & 时间段 & 充电量 & 放电量\\
\midrule
00:00--04:00 & 3924.13 & 181.53 & 04:00--08:00 & 1015.09 & 5240.21\\
08:00--12:00 & 6123.03 & 2777.38 & 12:00--16:00 & 3923.99 & 420.71\\
16:00--20:00 & 268.54 & 5719.49 & 20:00--24:00 & 5703.85 & 2960.43\\
\midrule
0:00储电量 & \multicolumn{2}{r}{6705.54} & 24:00储电量 & \multicolumn{2}{r}{6346.37}\\
\bottomrule
\end{tabularx}
\par\smallskip\textbf{2025-06-21}\par\smallskip
\begin{tabularx}{\linewidth}{lRRlRR}
\toprule
时间段 & 充电量 & 放电量 & 时间段 & 充电量 & 放电量\\
\midrule
00:00--04:00 & 1394.50 & 465.11 & 04:00--08:00 & 460.19 & 4390.39\\
08:00--12:00 & 6534.90 & 0.00 & 12:00--16:00 & 0.00 & 0.00\\
16:00--20:00 & 156.59 & 4882.20 & 20:00--24:00 & 6792.80 & 2485.90\\
\midrule
0:00储电量 & \multicolumn{2}{r}{8644.37} & 24:00储电量 & \multicolumn{2}{r}{8867.68}\\
\bottomrule
\end{tabularx}
\par\smallskip\textbf{2025-09-23}\par\smallskip
\begin{tabularx}{\linewidth}{lRRlRR}
\toprule
时间段 & 充电量 & 放电量 & 时间段 & 充电量 & 放电量\\
\midrule
00:00--04:00 & 4692.99 & 53.14 & 04:00--08:00 & 916.24 & 5893.48\\
08:00--12:00 & 5977.49 & 1896.62 & 12:00--16:00 & 3708.39 & 625.98\\
16:00--20:00 & 151.78 & 5739.35 & 20:00--24:00 & 5609.40 & 2640.89\\
\midrule
0:00储电量 & \multicolumn{2}{r}{6006.02} & 24:00储电量 & \multicolumn{2}{r}{6235.06}\\
\bottomrule
\end{tabularx}
\par\smallskip\textbf{2025-12-21}\par\smallskip
\begin{tabularx}{\linewidth}{lRRlRR}
\toprule
时间段 & 充电量 & 放电量 & 时间段 & 充电量 & 放电量\\
\midrule
00:00--04:00 & 4499.79 & 40.16 & 04:00--08:00 & 762.27 & 1359.20\\
08:00--12:00 & 5913.75 & 6490.57 & 12:00--16:00 & 5645.48 & 2130.65\\
16:00--20:00 & 0.00 & 5526.27 & 20:00--24:00 & 5771.81 & 2842.53\\
\midrule
0:00储电量 & \multicolumn{2}{r}{6794.81} & 24:00储电量 & \multicolumn{2}{r}{6695.97}\\
\bottomrule
\end{tabularx}
\end{table}

\begin{table}[htbp]\centering\small
\caption{问题2指定日期紧急购电（题目表3布局，单位：kWh）}
\label{tab:q2-emergency}
\setlength{\tabcolsep}{3pt}
\begin{tabularx}{\linewidth}{lRlRlRlR}\toprule
\multicolumn{2}{c}{2025-03-20} & \multicolumn{2}{c}{2025-06-21} & \multicolumn{2}{c}{2025-09-23} & \multicolumn{2}{c}{2025-12-21}\\
时间段 & 购电量 & 时间段 & 购电量 & 时间段 & 购电量 & 时间段 & 购电量\\\midrule
无 & 0.00 & 无 & 0.00 & 20:20--20:40 & 388.40 & 无 & 0.00\\
 &   &  &   & 20:50--21:10 & 90.05 &  &  \\
 &   &  &   & 21:20--21:50 & 34.20 &  &  \\
\midrule
合计 & 0.00 & 合计 & 0.00 & 合计 & 512.64 & 合计 & 0.00\\
\bottomrule\end{tabularx}\end{table}

四个指定日中，9月23日存在三段紧急购电，合计512.64 kWh，其余三个指定日为零。其常规计划费用43098.28元，紧急费用3285.98元，全天实际费用46384.26元；不能把全天费用只写成计划量乘价格。相邻正值时段才合并为同一紧急区间，间隔含零时保持分段。

表\ref{tab:q2-purchase}的指定时段及全天购电量均指零时计划，不含紧急量；后者单列于表\ref{tab:q2-emergency}。表\ref{tab:q2-battery}填写实际动作，各四小时块内充放电均可能为正，但每个10分钟段仍互斥。完整计划保存在\texttt{result2.xlsx}，日费用分解可由同源结果数据核对。
\FloatBarrier

\subsubsection{风险、备用及其联合作用}

\begin{table}[htbp]\centering\small
\caption{问题2风险与备用的完整对照}
\label{tab:v2_risk_modules}
\begin{tabularx}{\linewidth}{LRRRR}\toprule
策略 & 总费/万元 & 紧急量/kWh & 日费CVaR95/元 & 末库存/kWh\\\midrule
期望费用 & 1401.28 & 85017.69 & 70731.74 & 6240.32\\
固定风险 & 1411.84 & 56316.17 & 69028.62 & 6318.65\\
自适应风险 & 1405.41 & 66920.81 & 69476.71 & 6255.33\\
固定备用 & 1401.09 & 74987.90 & 70409.51 & 6240.32\\
动态备用 & 1402.30 & 76168.28 & 70299.39 & 6245.11\\
固定风险+备用 & 1413.36 & 49102.47 & 68862.93 & 6318.65\\
自适应风险+备用 & 1406.93 & 59081.98 & 69070.96 & 6260.12\\
自适应联合，12场景 & 1404.26 & 75336.25 & 70065.84 & 6289.57\\
\bottomrule\end{tabularx}
\par\smallskip\parbox{\linewidth}{\footnotesize 本表基于未启用候选筛选的母模型，各组保留原实验配置及数值，覆盖相同334天并连续推进状态；不作为最终筛选算法的同条件消融。}
\end{table}

\begin{table}[htbp]\centering\small
\caption{固定与自适应策略在评价期实现的平均保守程度}
\label{tab:v2_conservatism}
\begin{tabularx}{\linewidth}{LRR}\toprule
策略 & 平均风险权重 & 平均参考备用/kWh\\\midrule
固定风险 & 0.10368 & 0.00\\
自适应风险 & 0.05506 & 0.00\\
固定备用 & 0.00000 & 207.22\\
动态备用 & 0.00000 & 244.39\\
固定联合 & 0.10368 & 207.22\\
自适应联合 & 0.05506 & 244.39\\
\bottomrule\end{tabularx}
\par\smallskip\parbox{\linewidth}{\footnotesize 固定参数由一月动态序列的均值匹配；由于季节与误差分布变化，评价期均值不保证相等，不用评价期重新匹配。}
\end{table}

与主案相比，仅固定风险增加105631.25元，仅自适应风险增加41280.51元，紧急量分别降至56316.17和66920.81 kWh。固定备用单独使用节省1891.73元，动态备用单独使用增加10152.11元，说明备用的经济效果取决于其形状和时段作用。联合自适应策略增加56524.93元，紧急量降至59081.98 kWh。紧急费用CVaR的优化目标和表中日总费用CVaR95不是同一个损失对象，不据其中一项下降推断另一项必然下降。

固定参数虽由一月均值匹配，评价期固定和自适应的平均风险权重分别为0.103681与0.055059，参考备用均值分别为207.22与244.39 kWh，并未保持相同。因此两类方法之间的差异同时包含保守程度变化，不能全部归结为自适应对高风险时段识别得更准确。

以$C_{\rm joint}-C_{\rm risk}-C_{\rm reserve}+C_{\rm expectation}$描述联合作用，固定和自适应两组分别为17088.03元和5092.30元，均为正。这意味着本回测中联合费用高于两项单独费用变化的简单相加，但不是独立因果效应估计，也不是普遍的重复保守定理。

\begin{figure}[htbp]\centering
\includegraphics[width=0.95\linewidth]{figures/results/q2_risk_heatmap.pdf}
\caption{七种风险与备用设置相对主案的指标差异；前三行标注百分比，备用和末库存标注kWh差值，颜色仅在同一指标行内比较}
\label{fig:q2-risk-heatmap}
\end{figure}

图\ref{fig:q2-risk-heatmap}与完整对照表共同展示，紧急购电减少往往伴随额外常规承诺。费用、紧急量和日费CVaR95使用相对主案百分比变化；参考备用的主案值为零，故与末库存一样报告kWh差值。配色按各指标行的最大绝对变化归一，不跨行比较色深，也不将红色一律解释为不利。固定备用在评价期较便宜，也不据此替换预先冻结的主案。风险扩展回答的是为减少特定补购需求愿意支付多少费用，而非赛题额外规定了必须达到的安全阈值。

\subsubsection{28轨迹保留与12代表日压缩}

% Generated by code/make_question_tables.py; do not edit values.
% Source: data/experiments_v2.json; SHA256=e231ea14b234bd4593ec19c5f913ce028ed3b15ee677b4b58eafc7d3a658f29a
\begin{table}[htbp]\centering\small
\caption{问题二完整场景与压缩场景的同条件比较}
\label{tab:v2_scenarios}
\label{tab:v2_solve_times}
\begin{tabularx}{\linewidth}{LRRRR}\toprule
场景数 & 总费/万元 & 紧急量/kWh & 日费CVaR95/元 & 秒/日\\\midrule
28 & 1401.28 & 85017.69 & 70731.74 & 0.0521\\
12 & 1400.63 & 105407.24 & 72139.66 & 0.0276\\
\bottomrule\end{tabularx}
\par\smallskip\parbox{\linewidth}{\footnotesize 相同334天、预测、执行器及期初库存；末库存分别为6240.32和6257.26 kWh。时间包含预测、求解和执行，在本机并行条件下仅反映工程量级。}
\end{table}


经济目标下12场景比28场景少付6465.72元，却多购20389.55 kWh紧急电量，日费用CVaR95也更高。风险与备用联合开启时，12场景总费用14042623.21元、紧急量75336.25 kWh；28场景为14069332.59元、59081.98 kWh。压缩既改变期望费用，也会改变尾部及紧急需求，不能凭单一总费用接近就判断尾部已被保留。

28个等权历史日的10\%尾部仅有2.8条轨迹的概率质量，风险估计本来就受样本量限制；压缩后还可能集中到少数代表日。主案预定保留28条，排除额外压缩影响，求解耗时仍适合本题规模。该选择并不保证本年度实际账单最小，也不消除28日历史窗口有限性的限制。

为核对缩减程序的数值过程，另外记录334个日期窗口的实际medoid迭代。令标准化历史轨迹为$x_j$，第$m$轮代表集合为$\mathcal M^{(m)}$，压缩距离目标为
\begin{equation}
 J_m=\frac{1}{|\mathcal H_D|}\sum_{j\in\mathcal H_D}
       \min_{\ell\in\mathcal M^{(m)}}\|x_j-x_\ell\|_2.
 \label{eq:q2_compression_objective}
\end{equation}
图\ref{fig:q2-scenario-convergence}对每个窗口用其初值归一为$J_m/J_0$，再显示跨窗口中位数和P10--P90区间。334个窗口均达到代表稳定，迭代次数中位数2、最大5，最终相对目标中位数0.9816002。代表和权重已逐窗口与原缩减函数输出核对。

\begin{figure}[htbp]\centering
\includegraphics[width=0.93\linewidth]{figures/results/q2_scenario_convergence.pdf}
\caption{12代表日压缩距离目标的实际迭代；曲线为334窗口归一目标中位数，阴影为逐轮P10--P90}
\label{fig:q2-scenario-convergence}
\end{figure}

停止后的窗口只重复终值以对齐横轴，图中分位区间描述不同窗口之间的离散程度，不是置信带。此图证明的是历史轨迹几何压缩过程按预定准则稳定，并非LP求解收敛、经济费用下降或尾部风险保持；经济结论仍由表\ref{tab:v2_scenarios}及实际执行账单决定。
\FloatBarrier

\FloatBarrier
\section{问题三模型建立与求解}
\label{sec:question3}

\subsection{问题三分析}

问题三在问题二的日前购电基础上，增加0、6、12、18时发布的未来24小时光伏预报，并允许在规定时点修订尚未执行的常规购电计划。由此产生两个不同的决策频率：常规购电只能按交易时点修订，储能则每10分钟根据实际负荷、光伏与储电量动作。若把两类决策混为每10分钟重新购电，会绕过题目规定的调整费用；若完全照搬日前储能轨迹，又可能在预测误差出现后违反能量平衡。因此，本问采用外层采购、内层执行的双时间尺度结构。

需要先回答三个问题。第一，原计划已经付款，日内增加或减少购电分别意味着什么边际费用；第二，新预报和实测SOC各自能改善哪些决策；第三，规划模型使用的参考储能轨迹，与实际控制器产生的费用是否相同。本问据此保留受交易约束的情景LP作为母模型，在外层加入基于执行规则回放的购电候选筛选，储能仍按原冻结规则执行。最终方案在零时和三个日内修订时点均启用筛选，以开发期闭环账单及筛选时点消融检验其经济性，同时报告紧急补购增加的代价。

\subsection{问题三模型建立}

\subsubsection{结算口径、实际账单与下调支配性}
\label{sec:settlement_mechanism}

记$g_{D,t}$为日期$D$零时承诺的常规购电量，$a_{D,t}$为该时段最后生效的常规购电量，$h_{D,t}$为实际紧急购电量，三者单位均为kWh。设$[x]_+=\max\{x,0\}$。赛题给出计划量计费、1.5倍上调费用、50\%下调违约费及5倍紧急费用，但没有单列取消原计划的退款项。本文采用原计划全额付款且不退款的主解释，并保留退款解释作为敏感性对照。

\begin{table}[htbp]
\centering\small
\caption{问题三两种结算口径及单时段示例（不含紧急购电费）}
\label{tab:mech_settlement_convention}
\begin{tabularx}{\linewidth}{p{0.18\linewidth}Xrr}
\toprule
口径 & 非紧急费用 & 下调示例 & 上调示例\\
\midrule
主口径：不退款 & $pg+1.5p[a-g]_++0.5p[g-a]_+$ & 110元 & 130元\\
退款敏感性口径 & $p\min\{g,a\}+1.5p[a-g]_++0.5p[g-a]_+$ & 90元 & 130元\\
\bottomrule
\end{tabularx}
\par\smallskip\parbox{\linewidth}{\footnotesize 两例均取$g=100$ kWh、$p=1$元/kWh；下调时$a=80$ kWh，上调时$a=120$ kWh。退款解释仅返还取消部分的原计划费用，紧急费用在两种口径下均另加$5ph$。}
\end{table}

主口径的当日真实账单为
\begin{equation}
 C_D=\sum_{t=0}^{143}p_{D,t}\left(g_{D,t}+1.5[a_{D,t}-g_{D,t}]_+
 +0.5[g_{D,t}-a_{D,t}]_++5h_{D,t}\right).
 \label{eq:settlement}
\end{equation}
每次修订均以零时原计划为基准，按最后生效量结算，不把前后两次修订之间的差额重复收费。原计划费用、调整费用及紧急费用分别保存，辅助终端罚项不进入式\eqref{eq:settlement}。

在正电价、原计划不退款、富余电量允许免费弃置三个条件下，下调受到支配。若某可行解$a_t<g_t$，改为$a'_t=g_t$，保持充放电、紧急购电及SOC不变，将多出的$g_t-a_t$弃置，物理约束仍成立，费用却减少$0.5p_t(g_t-a_t)>0$。因而至少存在满足$a_t\ge g_t$的最优解。该结论不适用于退款解释，也不授权内层控制器自行改动已承诺合同。未启用筛选的母模型在两种口径下的334天费用相差320347.51元，故不能把结算语义当作次要实现细节，或因某种解释更便宜而改用该解释。

\subsubsection{由1.5倍增购费用推导条件采购阈值}

沿用问题二的无储能单时段分析，设$X$为修订时未知的实际净需求，$p>0$，并固定原计划$g\ge0$。在主口径下只需考虑$a\ge g$，已承诺的$pg$为常数，因此当前增购决策为
\begin{equation}
 J_3(a)=1.5p(a-g)+5p\,\mathbb E[(X-a)_+],\qquad a\ge g.
 \label{eq:mech_q3_cost}
\end{equation}
其中概率和期望均以该次修订已经获得的信息为条件。设$F_X(x)=\Pr(X\le x)$、$F_X(x^-)=\Pr(X<x)$，则
\begin{equation}
 \partial J_3(a)=p[5F_X(a^-)-3.5,\ 5F_X(a)-3.5].
 \label{eq:q3_adjustment_subgradient}
\end{equation}
内点$a>g$的最优条件为$F_X(a^-)\le0.7\le F_X(a)$，边界$a=g$的最优条件为$F_X(g)\ge0.7$。按下分位数定义，可取
\begin{equation}
 a_3^\star=\max\{g,F_X^{-1}(0.7)\},\qquad
 a_3^\star-g=[F_X^{-1}(0.7)-g]_+.
 \label{eq:mech_q3_quantile}
\end{equation}
离散场景和连续分布均满足上述次梯度刻画；在连续内点，一阶条件也可写为$1.5p-5p\Pr(X>a)=0$。70\%对应给定原计划后的最终有效购电量阈值，既不是增购量本身的分位数，也不是问题三零时计划的最优分位数。储能还引入跨时段库存价值，所以完整调度仍须联合优化，不能逐时套用该解析式。

\subsubsection{合法信息更新与外层滚动采购模型}

用十分钟物理索引表示外层时点，令$\mathcal T_u=\{0,36,72,108\}$。在$\tau\in\mathcal T_u$读取实测储电量$e_\tau^{\rm act}$，只处理$\tau,\ldots,143$的未执行区间；已执行动作及费用保持不变。光伏采用该时点刚发布的预报，负荷零时预测采用一月已选定的最近两个同星期日均值。为利用已经观察到的当天负荷偏差，在$\tau>0$时计算
\begin{align}
 b^L_{D,\tau}&=\frac{1}{12}\sum_{k=\tau-12}^{\tau-1}
       (L_{D,k}-\widehat L^{(0)}_{D,k}),\label{eq:q3_load_bias}\\
 \widehat L^{(\tau)}_{D,t}&=\max\left\{0,\widehat L^{(0)}_{D,t}+
 \operatorname{clip}\!\left(b^L_{D,\tau}e^{-(t-\tau+1)/36},
       -0.3\widehat L^{(0)}_{D,t},0.3\widehat L^{(0)}_{D,t}\right)\right\},\quad t\ge\tau.
 \label{eq:q3_load_update}
\end{align}
$\operatorname{clip}(x,l,u)=\min\{u,\max\{l,x\}\}$。该修正只使用此前两小时实测值，按六小时尺度衰减，并限制在零时预测的$\pm30\%$；这些是预先规定的有限反馈规则，没有用当日未来负荷校准。

在同一发布时间、相同剩余时域内提取此前最多28个完整日的净负荷残差，形成$n_{t,s}$及概率$\pi_s$。保留整日残差片段能表示连续偏高或偏低的误差。母模型在零时求解问题二式\eqref{eq:q2_objective}的经济场景模型；后续时点以已选定的零时计划$g_t$为基准，令$a_t-g_t=u_t^+-u_t^-$，$u_t^+,u_t^-\ge0$，求解
\begin{align}
 \min\quad &\sum_{t=\tau}^{143}\bar p_t(1.5u_t^++0.5u_t^-)
 +\sum_s\pi_s\sum_{t=\tau}^{143}5p_{t,s}h_{t,s}
 +\beta_T|e_{144}-6000|,\label{eq:q3_revision_objective}\\
 \text{s.t.}\quad&a_t+d_t+h_{t,s}-c_t\ge n_{t,s},\quad a_t,h_{t,s}\ge0,
                    \label{eq:q3_revision_balance}\\
 &a_t-g_t=u_t^+-u_t^-,\quad u_t^+,u_t^-\ge0,\quad
 e_\tau=e_\tau^{\rm act},\label{eq:q3_revision_contract}
\end{align}
并满足式\eqref{eq:soc}、\eqref{eq:limits}的储能约束。问题三$p_{t,s}=p_t$为已知分时电价，故$\bar p_t=p_t$；保留场景价格记法便于问题四沿用。原计划费用在当前修订中为常数，可从优化目标略去，但仍计入真实账单。终端系数$\beta_T=\operatorname{median}_{t\ge\tau}\bar p_t$，绝对值通过$e_{144}-6000=z^+-z^-$、$z^+,z^-\ge0$线性化。该项近似剩余库存价值，不强制每天回到6000 kWh。

外层参考充放电和SOC在场景间共享，紧急购电$h_{t,s}$作为场景补救量；主设置预先固定为28场景、无附加风险权重、无参考备用。每次求解当天全部剩余合同，但只把下一合法更新时间之前的部分推进为实际执行量，其后的未执行合同仍可在下个时点按同一规则修订。共享储能限制了策略类，不能把该外层目标直接称为实际控制器的期望费用。

\subsubsection{基于执行规则回放的购电候选筛选}
\label{sec:q3_candidate_screening}

为使采购决策的评估更接近实际执行，在母模型之外建立有限购电候选集合$\mathcal K_\tau$。零时包含五类候选：原情景LP、合法点预测LP，以及净需求预测分别增加$0.25\sigma_t$、$0.5\sigma_t$和$1.0\sigma_t$的三个误差余量LP。其中$\sigma_t^2=\sum_s\pi_s(n_{t,s}-\sum_j\pi_jn_{t,j})^2$，单位为kWh的平方；余量只增加点预测净需求，点预测价格保持不变。日内额外加入“保持当前剩余有效合同”候选，它保留此前修订后尚未执行的合同，并非恢复零时计划。各LP均使用当前同一实测SOC、储能物理约束及零风险、零备用设置，其目标值仅用于产生候选，不能直接跨候选比较。

记候选$k$的剩余常规购电量为$q_t^{(k)}$。对每个候选，从当前共同初态$e_\tau^{\rm act}$出发，在同一组历史情景上按式\eqref{eq:execute_charge}、\eqref{eq:execute_discharge}的固定规则逐段回放，获得$h_{t,s}^{(k)}$与$e_{144,s}^{(k)}$。回放满足$q_t^{(k)}+d_{t,s}^{(k)}+h_{t,s}^{(k)}-c_{t,s}^{(k)}\ge n_{t,s}$，允许富余弃置。所有购电和充放电变量均为电量，附件功率先乘$\Delta t=1/6$ h转换，费用不再重复乘时间步长。剩余时域账单定义为
\begin{align}
 C_{\tau,s}^{\mathrm{replay},(k)}
 &=\sum_{t=\tau}^{143}p_{t,s}\bigl(B_{\tau,t}^{(k)}+5h_{t,s}^{(k)}\bigr),
 \label{eq:q3_replay_bill}\\
 B_{\tau,t}^{(k)}&=
 \begin{cases}
 q_t^{(k)},&\tau=0,\\[1ex]
 g_t+1.5[q_t^{(k)}-g_t]_++0.5[g_t-q_t^{(k)}]_+,&\tau>0.
 \end{cases}\notag
\end{align}
其中$B_{\tau,t}^{(k)}$为以标准电价折算的非紧急计费量。因此，零时按候选计划量全额计费，日内仍计入原计划费，并将最后有效量相对零时计划的调整费完整计入；没有把购电费改成实际使用量的费用。已执行区间对本次选择为常数，不纳入剩余区间评分。

统一评分选择为
\begin{align}
 k_\tau^*&=\arg\min_{k\in\mathcal K_\tau}\sum_s\pi_s
 \left[C_{\tau,s}^{\mathrm{replay},(k)}
 +\beta_\tau\left|e_{144,s}^{(k)}-6000\right|\right],
 \label{eq:q3_candidate_score}\\
 \beta_\tau&=\operatorname{median}_{t\ge\tau}\sum_s\pi_sp_{t,s}.\notag
\end{align}
同一次筛选的情景、权重、初态、执行器、结算规则和终端系数对所有候选相同。末库存偏离项先逐情景计算绝对值再加权，仅用于候选评分，不计入式\eqref{eq:settlement}的实际费用。零时将所选$q^{(k_0^*)}$保存为$g$；日内更新剩余有效合同，并始终以本日$g$为调整基准。评分完全相同时按固定候选顺序选择，优先保留情景LP。

候选回放将当前合同固定至午夜，没有模拟后续尚未发布的预报与合同修订；实际运行仍在下一个合法时点读取新信息并重新决策。因此，这是一种有限候选的执行费用评估，部分改善了规划评估与实际控制的衔接，并非完整多阶段随机最优控制。扩大候选集只能保证同一情景评分下的最小估计费用不增，不能保证未来实际账单下降。最终方案在$\mathcal T_u$的四个时点均启用筛选，适用问题与时点由下文开发性消融确定。

\subsubsection{母模型执行器对照：固定合同内层LP}

母模型与最终增强方案均采用式\eqref{eq:execute_charge}、\eqref{eq:execute_discharge}的规则：富余时在物理上限内充电，缺口时在功率与可用电量允许范围内放电，再紧急补购。该规则仅需当期$a_t,n_t,e_t$，易于执行；但它不直接比较当前与未来电量价值，也不保证跟随外层参考SOC。以下内层LP及跨日续接是基于未启用候选筛选的母模型所作的执行器对照，其结果不作为最终增强方案的同条件消融。

当前时段$t$固定当天剩余常规合同$a_{k|t}$；令当前$\widetilde n_{t|t}=n_t$、$\widetilde p_{t|t}=p_t$，对$k>t$只使用最近合法预报，定义
\begin{equation}
 s_{k|t}=[a_{k|t}-\widetilde n_{k|t}]_+,\qquad
 b_{k|t}=[\widetilde n_{k|t}-a_{k|t}]_+.
 \label{eq:v2_surplus_deficit}
\end{equation}
内层优化为
\begin{align}
 \min\quad&\sum_{k=t}^{143}5\widetilde p_{k|t}(b_{k|t}-d_{k|t})
       +\beta_T|e_{144|t}-6000|,\label{eq:v2_inner_objective}\\
 \text{s.t.}\quad&0\le c_{k|t}\le\min\{U,s_{k|t}\},\quad
 0\le d_{k|t}\le\min\{U,b_{k|t}\},\label{eq:v2_inner_bounds}\\
 &h_{k|t}=b_{k|t}-d_{k|t},\quad
 w_{k|t}=s_{k|t}-c_{k|t},\quad e_{t|t}=e_t^{\rm act},
 \label{eq:v2_inner_balance}
\end{align}
同时满足SOC递推、容量和功率约束。$\beta_T$由当前剩余预测价格的中位数给定。常规合同费用此时不可改变，故从内层目标省略。只执行第一步动作，下一时段获得新实测量后再算；不能在内层增加或取消常规合同。

约束\eqref{eq:v2_inner_bounds}仅允许常规供电富余充电，紧急购电只补负载缺口，不用于主动给电池充电；富余与缺口不能同时为正，故充放电互斥自动成立。与规则执行不同，LP可在预测后续缺口更昂贵时减少当前放电，付出当期补购费以保留电量。作为行为核验，若电池可输出1 kWh、连续两段合同均为0且需求均为1 kWh、电价分别为1与2元/kWh，关闭终端罚项，则规则先放电、后付10元补购；LP先付5元补购、后放电。程序重现了这一机制，但该构造例子不代表全年必然节费。

\subsubsection{母模型执行器对照：跨日预报的续接价值}

6、12、18时公布的是未来24小时预报，分别含次日最初6、12、18小时，即$m=36,72,108$个十分钟段。跨日候选只扩展内层价值估计：次日光伏使用本次已经公布的预报，次日负荷由当日之前已完成的历史日外推；不将当日尚未发生的负荷作为历史样本。零时没有跨日部分，取$m=0$。

当$m>0$时，把固定日末罚项替换为
\begin{equation}
 \widehat V_D(e_{144})=\min_{z,c',d',e'}
 \left\{\sum_{j=0}^{m-1}\widehat p'_jz_j+\beta_T|e'_m-6000|\right\},
 \label{eq:v2_continuation}
\end{equation}
约束为$z_j\ge0$、$z_j+d'_j-c'_j\ge\widehat n'_j$、$e'_0=e_{144}$及相同储能状态方程和物理界。$z_j$为次日临时购电变量，$\beta_T$仍由当天剩余预测价格确定。续接费用使次日需求和补能条件能够影响当前保留库存的价值；这些次日变量不执行、不填入当天购电表、不计入当天真实账单，次日零时再制定正式计划。

该执行器候选仍为确定性续接近似，没有求解次日完整随机交易过程。外层零时模型及候选回放也尚未显式计入以后1.5倍修订的机会：事后允许重新优化，不等于日前已经计算了未来修订的信息价值。严格一致的多阶段策略还需描述合法发布时间上的后续合同修订与信息分支；不能给每条场景独立优化全天储能后直接执行。

\subsection{问题三求解}

\subsubsection{配置确定、滚动实施及费用统计}

先预定经济基准为28场景、风险权重和备用均关闭，再在1月15--31日按实际账单选择预测器与控制候选，不能将风险关闭误写为对风险参数进行一月搜索后的结论。所有同一问题的验证候选从相同实测初态出发，期末库存随候选列出。表\ref{tab:v2_controller_validation}给出未启用候选筛选的母模型在已选预测配置下的控制器验证：问题三规则执行费用919340.82元，剩余日LP为929987.84元，LP加跨日为920497.18元，因而冻结规则执行器。期末库存并不强制相等，这只是有限候选与有限验证期内的选择。预测器与实时执行器沿用这一冻结配置；候选筛选的适用问题及启用时点则依据2--12月开发性回测确定，最终配置确定后固定，本文所报告增强方案结果属于开发期闭环评价。

\begin{table}[htbp]\centering\small
\caption{相同已选预测配置下的执行器一月验证}
\label{tab:v2_controller_validation}
\begin{tabularx}{\linewidth}{LLRRR}\toprule
问题 & 执行器 & 实际费用/元 & 紧急量/kWh & 末库存/kWh\\\midrule
问题2 & 规则执行 & 917318.59 & 19.93 & 7155.31\\
问题2 & 储能LP & 928995.30 & 474.33 & 6000.00\\
问题3 & 规则执行 & 919340.82 & 49.23 & 7143.30\\
问题3 & 储能LP & 929987.84 & 195.13 & 6000.00\\
问题3 & 储能LP+跨日 & 920497.18 & 195.13 & 7143.30\\
问题4-2 & 规则执行 & 1042234.20 & 19.93 & 7298.41\\
问题4-2 & 储能LP & 1052652.76 & 19.93 & 6000.00\\
问题4-3 & 规则执行 & 1043044.13 & 44.10 & 7308.72\\
问题4-3 & 储能LP & 1057597.78 & 220.05 & 6000.00\\
问题4-3 & 储能LP+跨日 & 1047012.91 & 220.05 & 7291.46\\
\bottomrule\end{tabularx}
\par\smallskip\parbox{\linewidth}{\footnotesize 同一问题内使用相同1月15日期初库存，之后独立因果运行。费用不含终端罚项；末库存并不强制相等。}
\end{table}

实际求解使用稀疏LP和HiGHS。每日流程为：零时读取合法预报与承接SOC，生成购电候选并按式\eqref{eq:q3_candidate_score}筛选原计划；每10分钟读取当前实值并按固定规则执行储能；到6、12、18时用实测SOC和新信息生成并筛选未执行采购段，其中保留当前合同候选；日末按式\eqref{eq:settlement}分项结算，并令$e_{D+1,0}=e_{D,144}$。原版已执行的一月轨迹作为固定前缀，同一问题的母模型、最终方案及对照从相同2月初态开始，之后独立、连续推进状态，不能用后来选定的配置回写一月，也不每日重置库存。

问题三最终方案2月初态为8814.60 kWh，12月末为6240.32 kWh。334天计划费13503057.82元、调整费248699.70元、紧急费119713.93元，总费用13871471.44元。紧急购电21029.90 kWh，涉及119天、424个十分钟段；经紧急补购后逐段满足负载，紧急量不代表失供电量。实际日账单CVaR95为65552.76元，最高日费用为67941.61元。与保留母模型的问题二横向比较，总费用低1.0086\%、紧急量低75.26\%，但两问的零时光伏来源、初态和采购配置不同，不能把全部差异归为滚动反馈的独立作用。

\subsubsection{筛选时点消融与最终策略选择}
\label{sec:q3_screening_ablation}

为考察候选筛选在不同决策阶段的作用，在保持预测方法、历史情景构造、实时执行器、结算规则及合法采购修订时点不变的条件下，设置不启用筛选、仅零时筛选、仅日内筛选及零时与日内均筛选四种配置。关闭筛选的时点仍调用母模型情景LP，使用该方案自身的实测SOC与零时计划重新优化；“仅零时筛选”仍保留6、12、18时的合法预测更新和采购修订，并不等于仅制定一次日前计划。各组从同一问题的共同初态连续运行，以计划费用、调整费用和紧急费用之和评价，不将评分终端项计入实际账单。

\begin{table}[htbp]\centering\small
\caption{候选筛选启用时点对实际总费用的影响}
\label{tab:outer_timing_ablation}
\begin{tabularx}{\linewidth}{LRR}\toprule
候选筛选配置 & 问题3总费用/元 & 问题4滚动版总费用/元\\\midrule
不启用筛选：母模型基线 & 13,975,810.86 & 14,769,232.53\\
仅零时启用 & 13,911,957.70 & 14,712,596.45\\
仅日内启用 & 13,947,087.88 & 14,734,593.23\\
\textbf{零时＋日内均启用：最终方案} & \textbf{13,871,471.44} & \textbf{14,670,237.87}\\
\bottomrule\end{tabularx}
\par\smallskip\parbox{\linewidth}{\footnotesize 各组均保留0、6、12、18时的合法预报更新与采购修订，仅改变是否启用候选筛选；关闭筛选时仍调用原情景优化模型。各方案从同一问题的共同初态出发，之后独立、连续推进状态，不每日重置。评价范围为2025年2月1日至12月31日，共334天。}
\end{table}


在2025年2月1日至12月31日的334天开发性闭环回测中，全时点筛选在表\ref{tab:outer_timing_ablation}的四种配置中费用最低。第三问相较未启用筛选的母模型13975810.86元节省104339.41元，降幅0.747\%；相较仅零时和仅日内筛选，分别进一步节省40486.25元和75616.44元。第四问滚动版全时点费用14670237.87元，相较其母模型节省98994.66元，降幅0.670\%；相较仅零时和仅日内筛选，分别进一步节省42358.58元和64355.36元。差值由未舍入账单计算。这些结果支持在本次比较的四种配置中保留全时点筛选，不证明其在任意数据或候选集合下均最优。

\begin{table}[htbp]\centering\small
\caption{滚动采购候选筛选的费用与紧急补购权衡}
\label{tab:outer_screening_tradeoff}
\begin{tabularx}{\linewidth}{LLRRRR}\toprule
问题 & 配置 & 计划费/万元 & 调整费/万元 & 紧急费/万元 & 紧急量/kWh\\\midrule
问题3 & 母模型 & 1367.46 & 23.07 & 7.06 & 12647.76\\
问题3 & 最终方案 & 1350.31 & 24.87 & 11.97 & 21029.90\\
问题4滚动版 & 母模型 & 1446.02 & 24.19 & 6.71 & 11950.68\\
问题4滚动版 & 最终方案 & 1426.52 & 28.31 & 12.19 & 21293.41\\
\bottomrule\end{tabularx}
\par\smallskip\parbox{\linewidth}{\footnotesize 计划费按零时计划量全额结算，调整费始终以零时原计划为基准；紧急购电单独计费。终端偏离项不计入实际账单。两问题均以计划费用下降补偿新增调整费与紧急费。}
\end{table}


同一问题四组方案的评价期末库存一致，第三问均为6240.32 kWh、第四问滚动版均为6274.99 kWh，因此费用排序不能由期末额外消耗库存解释。但表\ref{tab:outer_screening_tradeoff}也显示，经济性改善伴随更多紧急补购：第三问紧急量由12647.76增至21029.90 kWh，第四问滚动版由11950.68增至21293.41 kWh。第三问计划费减少171508.13元，超过新增调整费18034.75元与紧急费49133.97元；第四问滚动版计划费减少195001.12元，超过新增调整费41200.07元与紧急费54806.39元。收益来自账单三项的净变化，并非所有运行指标同时改善。

同一候选模块用于问题二和第四问日前版时，费用反而分别增加22571.98元（0.161\%）和63356.15元（0.428\%），因此这两组仍采用母模型。最终分问配置为“日前两组不筛选、滚动两组零时与日内均筛选”。逐决策检查验证了当时可用信息边界、同口径评分和连续执行，但这种因果可执行性不等于模型选型没有参考回测结果；上述结论属于开发期选型证据，不是未接触数据上的独立测试，也不提供未来费用或跨年份最优性保证。

\subsubsection{状态、新光伏及负荷反馈的逐层贡献}

为分开信息来源，以下保留基于未启用候选筛选的母模型所作的独立消融。在相同初态、场景设定和规则控制器下，依次比较仅零时外层计划、只更新SOC并按零时预测重优化、再加入新光伏预报、最后加入已观测负荷偏差修正。四组均关闭候选筛选与跨日续接，避免“SOC-only”组从次日光伏预报间接获得新信息；每组之后按自身状态与合同因果演化。此实验考察信息来源，表\ref{tab:outer_timing_ablation}考察筛选模块的启用位置，两者配置与问题不同，不能相互替代。

\begin{table}[htbp]\centering\small
\caption{问题3母模型的状态、光伏和负荷反馈分解}
\label{tab:v2_feedback}
\begin{tabularx}{\linewidth}{LRRRR}\toprule
策略 & 总费/万元 & 紧急量/kWh & 日费CVaR95/元 & 末库存/kWh\\\midrule
仅零时外层计划 & 1416.90 & 80181.62 & 71801.23 & 6240.32\\
仅更新SOC并重优化 & 1399.17 & 39298.96 & 66748.14 & 6240.32\\
再加入新光伏预报 & 1407.57 & 40071.16 & 66732.97 & 6240.32\\
再加入负荷偏差反馈 & 1397.58 & 12647.76 & 65518.79 & 6240.32\\
\bottomrule\end{tabularx}
\par\smallskip\parbox{\linewidth}{\footnotesize 本表基于未启用候选筛选的母模型，各组保留原实验配置及数值，覆盖相同334天并连续推进状态；不作为最终筛选算法的同条件消融。}
\end{table}


仅零时外层计划费用14169011.89元，更新SOC后降至13991697.93元，减少177313.97元；再加入新光伏后增至14075665.67元，增加83967.74元；继续加入负荷偏差反馈后降至母模型基线13975810.86元，减少99854.81元。图\ref{fig:q3-feedback-waterfall}将这三项增量与起止账单分开显示。新信息的经济价值取决于采购、储能状态与误差的共同作用，不能预先认定每份新预报都节费。附件没有提供其他发布时间的预报，不用实际未来光伏代替不存在的高频预报来构造所谓实测收益。

\begin{figure}[htbp]\centering
\includegraphics[width=0.94\linewidth]{figures/results/q3_feedback_waterfall.pdf}
\caption{未启用候选筛选的母模型逐层反馈费用分解。横轴单位为万元，断轴用于放大账单附近的差异，起止条从零累计至334天总账单，中间条为相邻策略的费用差；增量包含各自闭环状态与后续采购的共同响应，不是单次预报的独立因果效应。}
\label{fig:q3-feedback-waterfall}
\end{figure}

\subsubsection{同计划回放：规划与控制的偏差}

本节分别考察原版12场景配置与未启用筛选的28场景母模型的规划与执行差异，诊断结果限于各自原有配置。原版采用12代表场景、自适应风险基数0.1和备用系数1，未启用候选筛选。重解全年3650个外层决策点，重建购电量在实际执行区间与保存结果完全相同。固定计划、初态与历史场景逐步回放，定义
\begin{equation}
 \Delta_{\rm exec}=\sum_s\pi_s\left(C_s^{\rm replay}-C_s^{\rm plan}\right).
 \label{eq:v2_execution_gap}
\end{equation}
两种费用都只含常规采购、调整及紧急采购，排除CVaR、终端罚项和数值破同解项。表\ref{tab:v2_execution_gap}汇总2--12月各决策点：原版问题二、三、四静态、四滚动的平均偏差依次为$-986.50$、$-569.90$、$-1068.47$、$-582.32$元。

\begin{table}[htbp]\centering\small
\caption{原版同计划、同初态、同场景的规划与执行偏差}
\label{tab:v2_execution_gap}
\begin{tabularx}{\linewidth}{LRRRR}\toprule
问题 & 决策点数 & 平均$\Delta_{\rm exec}$/元 & 负偏差比例 & SOC均方根偏差/kWh\\\midrule
问题2 & 334 & -986.50 & 97.90\% & 1319.72\\
问题3 & 1336 & -569.90 & 94.24\% & 994.32\\
问题4-2 & 334 & -1068.47 & 99.40\% & 1350.69\\
问题4-3 & 1336 & -582.32 & 93.94\% & 1048.30\\
\bottomrule\end{tabularx}
\par\smallskip\parbox{\linewidth}{\footnotesize 每一决策点先按场景概率计算费用差与SOC均方根偏差，再对决策点取均值；滚动剩余时域相互重叠，不将这些差求和解释为年度费用。全部费用仅计实际结算项目。}
\end{table}

负偏差不说明原规划已经正确描述控制，而是参考储能跨场景共享，实际规则可以利用当前观测改变动作，回放费用因此可能更低；正偏差也会出现。诊断窗口为剩余日，滚动窗口相互重叠，不能把差值相加成年度节费，也不能称其为相对全局最优的差距。图\ref{fig:q3-execution-band}取原版问题三正费用偏差最大的2025年2月9日零时决策，展示同一合同下12条历史场景的规则回放。此图来自保留个体轨迹的原版数据，与下文28场景母模型的16点双控制器诊断分开解释。

\begin{figure}[htbp]\centering
\includegraphics[width=0.95\linewidth]{figures/results/q3_execution_band.pdf}
\caption{原版问题三2025年2月9日零时同一合同下的12场景规则回放。虚线为共享参考SOC，实线为场景概率加权均值，细灰线为个体场景；阴影为各时刻按原概率计算的P10--P90加权经验分位带，不是置信区间或全年覆盖保证。图中不含内层LP轨迹。}
\label{fig:q3-execution-band}
\end{figure}

进一步在未启用筛选的28场景母模型上，对问题三4个指定日、4个外层时点，共16个决策点固定同一合同、实际初态与历史场景，对规则和内层LP各回放448条轨迹。每一步只揭示该步场景净需求和价格，未来仍使用该决策点已合法获得的点预测；后续合同和预报不再更新。这一计算分离了同一计划下换控制器的影响，以及全年闭环运行导致后续计划也改变的影响。

\begin{table}[htbp]\centering\small
\caption{未筛选母模型计划下两种执行器的回放诊断}
\label{tab:v2_inner_gap}
\begin{tabularx}{\linewidth}{LRR}\toprule
指标 & 规则执行 & 内层LP\\\midrule
有符号期望费用偏差/元 & -655.33 & -513.41\\
场景费用绝对偏差的期望/元 & 697.25 & 605.16\\
SOC均方根偏差/kWh & 1056.40 & 869.79\\
\bottomrule\end{tabularx}
\par\smallskip\parbox{\linewidth}{\footnotesize 四个指定日、四个时点共16个决策点，每点28场景。先计算场景概率加权指标，再对决策点等权平均；两种执行器固定相同购电合同和初态，仅逐步揭示当前场景值。}
\end{table}


平均SOC均方根偏差由规则执行1056.40 kWh降至内层LP的869.79 kWh，场景费用绝对偏差的期望由697.25元降至605.16元，但LP每个决策点的平均回放账单反而增加141.92元。16点中，10点的绝对期望费用偏差减小，仅1点的实际回放费用更低。更接近参考轨迹是一项可观察改善，却不能替代经济目标。最终增强方案由原执行器直接回放采购候选以缓解这一评估差异，但受剩余日固定合同近似所限，仍未达到严格多阶段一致。

\subsubsection{母模型内层LP与跨日续接的全年检验}

本节各组均未启用候选筛选，保持原有执行器对照配置。

\begin{table}[htbp]\centering\small
\caption{问题3母模型的执行控制与跨日续接对照}
\label{tab:v2_controllers}
\begin{tabularx}{\linewidth}{LRRRR}\toprule
策略 & 总费/万元 & 紧急量/kWh & 日费CVaR95/元 & 末库存/kWh\\\midrule
规则执行（母模型） & 1397.58 & 12647.76 & 65518.79 & 6240.32\\
剩余日储能LP & 1421.32 & 23825.27 & 66397.37 & 5934.55\\
储能LP+跨日价值 & 1405.26 & 22407.53 & 66033.05 & 6240.32\\
\bottomrule\end{tabularx}
\par\smallskip\parbox{\linewidth}{\footnotesize 本表基于未启用候选筛选的母模型，各组保留原实验配置及数值，覆盖相同334天并连续推进状态；不作为最终筛选算法的同条件消融。}
\end{table}


在334天完整闭环中，仅将母模型规则执行器换成剩余日LP，费用为14213193.28元，比规则母模型增加237382.42元，紧急量也增至23825.27 kWh。该候选将外层共享参考储能与内层有限点预测组合使用，并不必然优化真实未来费用，因此不予采用。

在相同LP下启用合法跨日信息后，费用降至14052638.61元，比截断于午夜的LP少160554.67元，紧急量少1417.74 kWh，年末库存还增加305.77 kWh。跨日改善并非依靠更多放空电池取得，但它仍比规则母模型多付76827.75元。控制器和续接候选均保持一月后的设置冻结，不依据评价期费用重新选择执行器；外层候选筛选的开发期选型另见第\ref{sec:q3_screening_ablation}节。

\subsubsection{退款口径的独立敏感性}

固定未启用候选筛选的母模型的预测、场景、控制器和初态，仅切换表\ref{tab:mech_settlement_convention}中的退款解释并重新优化，按各自规则结算。该比较检验制度语义，既不改变附件主口径，也不计作算法节费；表中不退款基线仍为13975810.86元，不与最终筛选方案混列。

\begin{table}[htbp]\centering\small
\caption{问题3母模型的结算解释敏感性}
\label{tab:v2_settlement}
\begin{tabularx}{\linewidth}{LRRRR}\toprule
策略 & 总费/万元 & 紧急量/kWh & 日费CVaR95/元 & 末库存/kWh\\\midrule
主口径：不退款 & 1397.58 & 12647.76 & 65518.79 & 6240.32\\
取消部分退款 & 1365.55 & 12871.85 & 64278.33 & 6240.32\\
\bottomrule\end{tabularx}
\par\smallskip\parbox{\linewidth}{\footnotesize 本表基于未启用候选筛选的母模型，各组保留原实验配置及数值，覆盖相同334天并连续推进状态；不作为最终筛选算法的同条件消融。}
\end{table}


退款解释费用13655463.35元，比不退款主口径少320347.51元，累计下调1096916.99 kWh，净调整费用为$-113349.02$元。两组年末库存相同，差异来自费用规则和对应采购响应，其量级大于若干模型增强的收益。退款后取消采购的边际费用改变，下调可以为正；不能再使用不退款条件下的支配结论。

\subsubsection{指定日期结果与实际运行轨迹}

图\ref{fig:v2-control}取最终方案评价期紧急量最大日7月4日，将最终上调量、实际SOC及紧急量分开呈现；库存子图同时给出未启用筛选的母模型实际SOC，便于观察不同采购决策形成的连续状态轨迹。两条曲线均来自各自实际执行数组。图\ref{fig:v2-adjustments}另列最终方案四个指定日的$a_t-g_t$，避免两条大部分重合的购电总量曲线遮盖修订。

\begin{figure}[htbp]\centering
\includegraphics[width=0.95\linewidth]{figures/results/v2_control_difference.pdf}
\caption{问题三最终方案紧急量最大日的购电调整、实际SOC与紧急购电。库存子图比较最终方案与未筛选母模型各自的实际SOC；虚线标示6、12、18时的日内修订时点，实际状态连续传递。}
\label{fig:v2-control}
\end{figure}
\begin{figure}[htbp]\centering
\includegraphics[width=0.95\linewidth]{figures/results/v2_adjustments.pdf}
\caption{最终方案四个指定日期最后生效常规购电量与零时计划之差$a_t-g_t$。数据来自全时点筛选的实际执行数组，不含紧急购电量；无调整时差值为零。}
\label{fig:v2-adjustments}
\end{figure}

\begin{table}[htbp]\centering\small
\caption{问题3指定日期购电结果（沿用题目表1布局，电量单位：kWh）}
\label{tab:q3-purchase}
\par\smallskip\textbf{2025-03-20}\par\smallskip
\begin{tabularx}{\linewidth}{lRlRlR}
\toprule
时间段 & 购电量 & 时间段 & 购电量 & 时间段 & 购电量\\
\midrule
10:00--10:10 & 0.00 & 12:00--12:10 & 624.29 & 14:00--14:10 & 0.00\\
16:00--16:10 & 671.04 & 18:00--18:10 & 704.20 & 20:00--20:10 & 0.00\\
\midrule
全天购电量 & 71789.01 & \multicolumn{2}{c}{全天购电费（元）} & \multicolumn{2}{r}{45563.84}\\
\bottomrule
\end{tabularx}
\par\smallskip\textbf{2025-06-21}\par\smallskip
\begin{tabularx}{\linewidth}{lRlRlR}
\toprule
时间段 & 购电量 & 时间段 & 购电量 & 时间段 & 购电量\\
\midrule
10:00--10:10 & 0.00 & 12:00--12:10 & 0.00 & 14:00--14:10 & 0.00\\
16:00--16:10 & 260.52 & 18:00--18:10 & 0.00 & 20:00--20:10 & 0.00\\
\midrule
全天购电量 & 37273.13 & \multicolumn{2}{c}{全天购电费（元）} & \multicolumn{2}{r}{21986.55}\\
\bottomrule
\end{tabularx}
\par\smallskip\textbf{2025-09-23}\par\smallskip
\begin{tabularx}{\linewidth}{lRlRlR}
\toprule
时间段 & 购电量 & 时间段 & 购电量 & 时间段 & 购电量\\
\midrule
10:00--10:10 & 0.00 & 12:00--12:10 & 444.81 & 14:00--14:10 & 0.00\\
16:00--16:10 & 611.01 & 18:00--18:10 & 801.37 & 20:00--20:10 & 0.00\\
\midrule
全天购电量 & 73472.67 & \multicolumn{2}{c}{全天购电费（元）} & \multicolumn{2}{r}{46201.57}\\
\bottomrule
\end{tabularx}
\par\smallskip\textbf{2025-12-21}\par\smallskip
\begin{tabularx}{\linewidth}{lRlRlR}
\toprule
时间段 & 购电量 & 时间段 & 购电量 & 时间段 & 购电量\\
\midrule
10:00--10:10 & 0.00 & 12:00--12:10 & 1054.19 & 14:00--14:10 & 0.00\\
16:00--16:10 & 877.48 & 18:00--18:10 & 744.23 & 20:00--20:10 & 0.00\\
\midrule
全天购电量 & 101747.42 & \multicolumn{2}{c}{全天购电费（元）} & \multicolumn{2}{r}{65702.81}\\
\bottomrule
\end{tabularx}
\par\smallskip\parbox{\linewidth}{\footnotesize 注：问题2与4-2列示计划购电量；问题3与4-3列示最终生效的调整后常规购电量。全天购电量不含紧急购电，全天购电费含计划、调整与紧急费用。}
\end{table}

\begin{table}[htbp]\centering\small
\caption{问题3指定日期储能结果（沿用题目表2布局，单位：kWh）}
\label{tab:q3-battery}
\par\smallskip\textbf{2025-03-20}\par\smallskip
\begin{tabularx}{\linewidth}{lRRlRR}
\toprule
时间段 & 充电量 & 放电量 & 时间段 & 充电量 & 放电量\\
\midrule
00:00--04:00 & 4146.19 & 181.53 & 04:00--08:00 & 883.67 & 6201.85\\
08:00--12:00 & 5765.88 & 2465.92 & 12:00--16:00 & 5215.25 & 254.72\\
16:00--20:00 & 1.31 & 5594.70 & 20:00--24:00 & 5703.85 & 2960.43\\
\midrule
0:00储电量 & \multicolumn{2}{r}{6505.69} & 24:00储电量 & \multicolumn{2}{r}{6428.95}\\
\bottomrule
\end{tabularx}
\par\smallskip\textbf{2025-06-21}\par\smallskip
\begin{tabularx}{\linewidth}{lRRlRR}
\toprule
时间段 & 充电量 & 放电量 & 时间段 & 充电量 & 放电量\\
\midrule
00:00--04:00 & 915.21 & 972.82 & 04:00--08:00 & 511.20 & 4390.39\\
08:00--12:00 & 8439.30 & 0.00 & 12:00--16:00 & 0.00 & 0.00\\
16:00--20:00 & 241.25 & 5157.72 & 20:00--24:00 & 6354.48 & 2485.90\\
\midrule
0:00储电量 & \multicolumn{2}{r}{7879.99} & 24:00储电量 & \multicolumn{2}{r}{8243.25}\\
\bottomrule
\end{tabularx}
\par\smallskip\textbf{2025-09-23}\par\smallskip
\begin{tabularx}{\linewidth}{lRRlRR}
\toprule
时间段 & 充电量 & 放电量 & 时间段 & 充电量 & 放电量\\
\midrule
00:00--04:00 & 4692.99 & 53.14 & 04:00--08:00 & 892.10 & 5612.58\\
08:00--12:00 & 6293.08 & 1896.62 & 12:00--16:00 & 3282.98 & 403.62\\
16:00--20:00 & 15.07 & 5619.05 & 20:00--24:00 & 5609.40 & 3045.21\\
\midrule
0:00储电量 & \multicolumn{2}{r}{6006.02} & 24:00储电量 & \multicolumn{2}{r}{6235.06}\\
\bottomrule
\end{tabularx}
\par\smallskip\textbf{2025-12-21}\par\smallskip
\begin{tabularx}{\linewidth}{lRRlRR}
\toprule
时间段 & 充电量 & 放电量 & 时间段 & 充电量 & 放电量\\
\midrule
00:00--04:00 & 4496.41 & 40.16 & 04:00--08:00 & 762.28 & 1406.28\\
08:00--12:00 & 6441.76 & 6526.43 & 12:00--16:00 & 5333.82 & 2222.97\\
16:00--20:00 & 0.00 & 5526.27 & 20:00--24:00 & 5771.81 & 2842.53\\
\midrule
0:00储电量 & \multicolumn{2}{r}{6797.85} & 24:00储电量 & \multicolumn{2}{r}{6695.97}\\
\bottomrule
\end{tabularx}
\end{table}

\begin{table}[htbp]\centering\small
\caption{问题3指定日期紧急购电（题目表3布局，单位：kWh）}
\label{tab:q3-emergency}
\setlength{\tabcolsep}{3pt}
\begin{tabularx}{\linewidth}{lRlRlRlR}\toprule
\multicolumn{2}{c}{2025-03-20} & \multicolumn{2}{c}{2025-06-21} & \multicolumn{2}{c}{2025-09-23} & \multicolumn{2}{c}{2025-12-21}\\
时间段 & 购电量 & 时间段 & 购电量 & 时间段 & 购电量 & 时间段 & 购电量\\\midrule
09:10--10:00 & 311.46 & 无 & 0.00 & 20:20--20:30 & 6.38 & 无 & 0.00\\
 &   &  &   & 20:50--21:10 & 67.74 &  &  \\
 &   &  &   & 21:20--21:50 & 34.20 &  &  \\
\midrule
合计 & 311.46 & 合计 & 0.00 & 合计 & 108.31 & 合计 & 0.00\\
\bottomrule\end{tabularx}\end{table}
\begin{table}[htbp]\centering\small
\caption{问题3指定日期计划、调整与费用分解}
\begin{tabularx}{\linewidth}{lRRRRR}\toprule
日期 & 计划量/kWh & 生效量/kWh & 计划费/元 & 调整费/元 & 紧急费/元\\\midrule
2025-03-20 & 71671.98 & 71789.01 & 43906.14 & 136.18 & 1521.51\\
2025-06-21 & 37259.86 & 37273.13 & 21962.65 & 23.90 & 0.00\\
2025-09-23 & 73472.67 & 73472.67 & 45684.65 & 0.00 & 516.92\\
2025-12-21 & 101747.42 & 101747.42 & 65702.81 & 0.00 & 0.00\\
\bottomrule\end{tabularx}\end{table}

以上指定日期表、运行轨迹及全年工作簿均由最终全时点筛选的同一执行数组生成。指定时段和紧急区间按统一物理索引生成，四小时充放电量按其中24个十分钟段合计。一个四小时区间内充、放电量同时为正，可以表示不同十分钟时段分别充电和放电。计划页保存原计划费用，调整页保存最后生效的常规量及完整实际费用；不把辅助目标值填成实际账单。
\FloatBarrier

\FloatBarrier
\section{问题四模型建立与求解}
\label{sec:question4}

\subsection{问题四分析}

问题四将已知分时电价改为随时间变化的交易电价，并要求分别推广问题二的日前采购和问题三的滚动策略。此时采购数量不仅受净负荷预测影响，还取决于缺口出现时的价格：同样少买1 kWh，若恰逢高价，五倍紧急补购造成的损失更大。因此，将负荷误差与价格误差分别独立抽样，可能破坏与费用直接相关的联合结构。

本问保留电池的状态方程、容量与功率限制，以及问题三已经确定的结算口径，改变外层的不确定输入和相应经济系数。附件4给出的实际价格用于逐时结算，未来尚未揭示的实际价格不能进入日前或滚动决策。先用价格加权分位数解释联合建模的原因，再以同日残差构造价荷场景；日前版保留母模型，滚动版沿用第\ref{sec:q3_candidate_screening}节的全时点候选筛选，最终比较两种可执行策略的实际费用。两案的预测配置、候选筛选设置与一月形成的初态不同，横向差异只作整体策略比较，不能全归因于单一模块。

\subsection{问题四模型建立}

\subsubsection{价格与需求相关时的采购分位数}

先考虑无储能单时段子问题。设未来净需求$X$与未来结算电价$P$可能相关，日前在已有信息下固定$q\ge0$；假设$P>0$几乎处处成立、$0<\mathbb E[P]<\infty$、$\mathbb E[PX_+]<\infty$，且富余购电无残值。期望费用为
\begin{equation}
 J_4(q)=\mathbb E[P]q+5\mathbb E[P(X-q)_+],\qquad q\ge0.
 \label{eq:mech_q4_cost}
\end{equation}
若将价格直接用均值替代，通常会遗漏$P$与缺口$(X-q)_+$的关联。为保留这一作用，定义
\begin{equation}
 G(x)=\frac{\mathbb E[P\,\mathbf 1_{\{X\le x\}}]}{\mathbb E[P]},\qquad
 G(x^-)=\frac{\mathbb E[P\,\mathbf 1_{\{X<x\}}]}{\mathbb E[P]}.
 \label{eq:mech_weighted_cdf}
\end{equation}
正价与有限期望保证$G$是合法分布函数。式\eqref{eq:mech_q4_cost}为凸函数，其次梯度为
\begin{equation}
 \partial J_4(q)=\mathbb E[P][5G(q^-)-4,\ 5G(q)-4].
 \label{eq:mech_q4_subgradient}
\end{equation}
内点最优条件为$G(q^-)\le0.8\le G(q)$，若最优点处连续，则
$\mathbb E[P\mathbf 1_{\{X\le q\}}]/\mathbb E[P]=0.8$。结合非负约束，可取
\begin{equation}
 q_4^\star=\max\{0,G^{-1}(0.8)\}.
 \label{eq:mech_weighted_quantile}
\end{equation}
对于离散场景$(x_s,p_s,\pi_s)$，只需把概率改写为
\begin{equation}
 \widetilde\pi_s=\frac{\pi_sp_s}{\sum_j\pi_jp_j},\qquad
 G(x)=\sum_{s:x_s\le x}\widetilde\pi_s.
 \label{eq:mech_scenario_reweight}
\end{equation}
于是最优采购对应价格加权的需求80\%分位数。若价荷独立，$G=F_X$，退化为问题二的普通分位数；若高需求更常伴随高价，该类场景的经济权重上升。该推导解释为何要保留联合关系，但只适用于无储能单时段。完整模型中的容量、效率、功率及跨时段能量价值，仍须由联合优化处理。

\subsubsection{合法价格预测与同日价荷残差场景}

以问题二、三已经选定的负荷和光伏预测配置为基础，只在1月验证期比较候选价格预测器，再冻结用于2--12月。动态日前案采用最近四个同星期日价格均值；动态滚动案采用最近两个同星期日均值。日前案负荷为最近两个同星期日均值、光伏为近三日均值；滚动案使用附件发布的光伏预报，并保留问题三的负荷偏差反馈。

动态滚动案在$\tau>0$时对价格采用与式\eqref{eq:q3_load_bias}、\eqref{eq:q3_load_update}相同的因果偏差修正：用此前12段实际价格相对零时预测的均值偏差，按六小时尺度衰减，并限制在零时预测的$\pm30\%$；价格预测正值下界取$10^{-6}$元/kWh。这些计算仅使用已经观察到的价格，不读取当天剩余实际价格。

对过去完整日$j<D$，分别求同发布时间的净负荷残差$\epsilon^n_{j,t}$和价格残差
\begin{equation}
 \epsilon^p_{j,t}=p_{j,t}-\widehat p^{(\tau)}_{j,t},\qquad t\ge\tau.
 \label{eq:q4_price_residual}
\end{equation}
两类残差按同一历史日配对，构成
\begin{equation}
 n_{t,s}=\widehat n_t+\epsilon^n_{j_s,t},\qquad
 p_{t,s}=\max\{0.01,\widehat p_t+\epsilon^p_{j_s,t}\}.
 \label{eq:joint_scenario}
\end{equation}
历史窗口至多28日，主案完整保留、等权赋值；价荷不拆日、不逐点重排。这样同时保存日内连续误差和价荷联系。若在12场景对照中进行代表日压缩，则先分别标准化净负荷与价格残差块再计算轨迹距离，防止不同单位的数值尺度支配选择。

式\eqref{eq:joint_scenario}的0.01元/kWh仅是外层加性场景的正价下界，并未改写实际价格序列。原始最低价0.0076元/kWh仍用于实际账单；实际结算价、合法点预测及合成场景价是三个不同对象。图\ref{fig:q4-correlation}在策略冻结后，对2--12月48096个时段的负荷、光伏、净负荷与电价作描述性Pearson相关分析。净负荷与价格的实际值相关系数为0.2620，按当时合法预测计算的残差相关系数为0.3468，显示分开两类残差会改变样本中的联合关系。该图不参与一月选型或未来场景构造，不授权提前使用评价期真值；未作时序依赖校正或显著性推断，相关系数也不表示因果效应。

\begin{figure}[htbp]\centering
\includegraphics[width=0.91\linewidth]{figures/results/q4_correlation_heatmap.pdf}
\caption{问题四滚动配置下48096个评价时段的Pearson相关矩阵，变量依次为负荷、光伏、净负荷、电价。左图为实际值，右图为实际值减各时段最近合法预测的误差，按0、6、12、18时各覆盖36段拼接；色标为$-1$至1。矩阵仅作事后描述，不用于选型或构造当天未来场景，未作时序相关校正或因果推断。}
\label{fig:q4-correlation}
\end{figure}

\subsubsection{动态日前采购与动态滚动采购的共同优化结构}

在每个外层时点，记$\bar p_t=\sum_s\pi_sp_{t,s}$，共享参考储能仍为$c_t,d_t,e_{t+1}$，紧急补购为$h_{t,s}$。动态日前案在零时求解
\begin{equation}
 \min\left\{\sum_{t=0}^{143}\bar p_tg_t+
  \sum_s\pi_s\sum_{t=0}^{143}5p_{t,s}h_{t,s}
  +\beta_T|e_{144}-6000|\right\},
 \label{eq:q4_day_ahead_objective}
\end{equation}
满足式\eqref{eq:scenario_balance}及储能物理约束，初态取当日实际承接库存。该模型的常规采购按场景均价评价，紧急费用则逐场景乘对应价格，因此保留了高缺口与高价同时出现时的损失。

动态滚动母模型的零时模型同式\eqref{eq:q4_day_ahead_objective}，在6、12、18时采用式\eqref{eq:q3_revision_objective}--\eqref{eq:q3_revision_contract}，把固定电价替换为本次联合场景价格。最终滚动案在四个合法时点均将该LP作为一个购电候选，与点预测、三个误差余量以及日内当前合同候选按式\eqref{eq:q3_candidate_score}统一筛选。每次回放完整保留式\eqref{eq:joint_scenario}中同一历史日的价格--净负荷配对，情景$s$的紧急费用使用$5p_{t,s}h_{t,s}^{(k)}$；不能把所有情景的紧急量统一乘平均价格。常规计划费和调整费同样在各价格情景下计算，终端系数仍为情景加权价格序列的中位数，各候选共用该系数。

其当日真实费用仍按式\eqref{eq:settlement}计算，式中的$p_{D,t}$此时取附件4已经实现的交易价格。日前案$a_t=g_t$、没有调整费，滚动案仅在合法时点更新合同，原计划始终按计划量全额结算；两者均由当前实际供需驱动储能执行，并连续传递日末SOC。回放合同固定至午夜的近似边界同问题三，不包含未来尚未发布的预报或以后可能发生的合同修订。

主案继续预先设置$\lambda=0$、$\rho_t=0$。若采用内层LP候选，则当前价格使用已观测实价、未来价格使用合法预测，合同固定规则与跨日临时变量含义均沿用问题三。该扩展没有将场景独立的全天储能计划作为可执行控制，也没有因为电价波动就额外假设可以随时无成本修改采购。

\subsection{问题四求解}

\subsubsection{最终配置与实际求解步骤}

一月验证中，在已选预测配置下，动态日前案规则执行费用1042234.20元，LP候选为1052652.76元；动态滚动案规则执行为1043044.13元，剩余日LP为1057597.78元，LP加跨日为1047012.91元，详见表\ref{tab:v2_controller_validation}。该比较基于未启用候选筛选的母模型，两案据此冻结规则执行器，不根据2--12月费用重新选择执行器。零风险与零备用仍是预定经济基准，并非从验证数据中搜索出来的参数最优结论。采购筛选的适用范围与时点另依据开发期回测确定：日前版启用筛选增费0.428\%，故保留母模型；滚动版采用表\ref{tab:outer_timing_ablation}中费用最低的零时与日内均筛选配置。最终配置固定后的报告口径为开发期闭环评价。

求解时先读取该时点已发布预报和已实现历史，构造同日价荷残差场景；日前版由情景LP直接确定常规计划，滚动版由各候选LP和同规则回放确定所选合同；按当前实际净负荷执行储能、记录必要紧急补购；最后按实际价格逐项计算账单。所有未来真实负荷、光伏和价格仅在其实现后进入执行与结算，场景均价不会代替实际价格填写费用。日前案与滚动案的2月初始SOC分别为9251.54、9030.40 kWh，均来自各自固定的一月运行前缀。

\subsubsection{全年费用、紧急需求与价格不确定性的影响}

动态日前案334天计划费用14280142.03元、紧急费用513982.26元，总费用14794124.29元，紧急量80544.00 kWh，涉及106天、596个十分钟时段。最终动态滚动案计划费用14265228.71元、调整费用283142.26元、紧急费用121866.90元，总费用14670237.87元，紧急量21293.41 kWh，涉及116天、418个十分钟时段。两案总费用相差123886.42元，即滚动案低0.8374\%，紧急量低73.56\%。

图\ref{fig:q4-cost-waterfall}把日前母模型到最终滚动案的总费用变化拆成计划、调整及紧急三项：滚动案计划费用减少14913.32元、紧急费用减少392115.36元，同时新增调整费用283142.26元。三项使用未舍入数据计算，总差以完整账单为准。减少昂贵补购不必然带来同比例的总节费，额外承诺和修订仍要付款。

\begin{figure}[htbp]\centering
\includegraphics[width=0.95\linewidth]{figures/results/q4_cost_waterfall.pdf}
\caption{动态日前母模型到全时点筛选滚动案的实际费用分解。横轴单位为万元并明确断轴，起止条从零累计至334天总费用，中间条分别为计划费、调整费和紧急费的变化，均来自最终配置的执行结果源；两案预测配置、筛选设置与初态不同，因此这是整体策略账单差，不是单一模块的因果节费。}
\label{fig:q4-cost-waterfall}
\end{figure}

最终动态滚动案实际日总费用CVaR95为77575.87元，动态日前案为83044.31元；最高日费用分别为81797.18元和108080.37元。这里的CVaR根据334个已实现日账单计算，是描述性尾部指标，不是启用额外风险惩罚的效果证明。两案期末库存分别为6274.99与6269.86 kWh，库存差与预测、初态差异一并报告；不把账单比较称为全局最优策略间的差距。

针对筛选模块本身，应与本问未筛选的滚动母模型14769232.53元比较，而非用日前案代替同条件基线。表\ref{tab:outer_timing_ablation}与表\ref{tab:outer_screening_tradeoff}表明，全时点筛选比该基线节省98994.66元，但紧急量增加9342.73 kWh，日费用CVaR95也由77420.84元升至77575.87元。因而本次选择体现为平均账单经济性与紧急补购、尾部费用之间的权衡，不能描述为各项指标全面改善。

\subsubsection{指定日期采购、储能与紧急购电结果}

以下分别给出动态日前母模型和最终全时点筛选滚动案在四个指定日期的六个指定时段采购量、全天常规量、储能四小时汇总及紧急购电区间。购电量均为kWh，完整费用包含紧急补购；计划费用和调整费用另在分解表中列出。表格、运行轨迹及结果工作簿均由各自最终配置的执行数组生成，使用同一物理时间索引，保留原模板显示标签，不再按标签循环平移数据。

\begin{table}[htbp]\centering\small
\caption{问题4-2指定日期购电结果（沿用题目表1布局，电量单位：kWh）}
\label{tab:q4_2-purchase}
\par\smallskip\textbf{2025-03-20}\par\smallskip
\begin{tabularx}{\linewidth}{lRlRlR}
\toprule
时间段 & 购电量 & 时间段 & 购电量 & 时间段 & 购电量\\
\midrule
10:00--10:10 & 0.00 & 12:00--12:10 & 564.20 & 14:00--14:10 & 0.00\\
16:00--16:10 & 490.32 & 18:00--18:10 & 0.00 & 20:00--20:10 & 743.27\\
\midrule
全天购电量 & 67869.67 & \multicolumn{2}{c}{全天购电费（元）} & \multicolumn{2}{r}{42207.16}\\
\bottomrule
\end{tabularx}
\par\smallskip\textbf{2025-06-21}\par\smallskip
\begin{tabularx}{\linewidth}{lRlRlR}
\toprule
时间段 & 购电量 & 时间段 & 购电量 & 时间段 & 购电量\\
\midrule
10:00--10:10 & 0.00 & 12:00--12:10 & 0.00 & 14:00--14:10 & 0.00\\
16:00--16:10 & 296.43 & 18:00--18:10 & 510.07 & 20:00--20:10 & 0.00\\
\midrule
全天购电量 & 40675.72 & \multicolumn{2}{c}{全天购电费（元）} & \multicolumn{2}{r}{21012.94}\\
\bottomrule
\end{tabularx}
\par\smallskip\textbf{2025-09-23}\par\smallskip
\begin{tabularx}{\linewidth}{lRlRlR}
\toprule
时间段 & 购电量 & 时间段 & 购电量 & 时间段 & 购电量\\
\midrule
10:00--10:10 & 0.00 & 12:00--12:10 & 497.24 & 14:00--14:10 & 0.00\\
16:00--16:10 & 537.29 & 18:00--18:10 & 0.00 & 20:00--20:10 & 0.00\\
\midrule
全天购电量 & 69762.68 & \multicolumn{2}{c}{全天购电费（元）} & \multicolumn{2}{r}{48309.88}\\
\bottomrule
\end{tabularx}
\par\smallskip\textbf{2025-12-21}\par\smallskip
\begin{tabularx}{\linewidth}{lRlRlR}
\toprule
时间段 & 购电量 & 时间段 & 购电量 & 时间段 & 购电量\\
\midrule
10:00--10:10 & 0.00 & 12:00--12:10 & 970.92 & 14:00--14:10 & 0.00\\
16:00--16:10 & 851.12 & 18:00--18:10 & 744.23 & 20:00--20:10 & 0.00\\
\midrule
全天购电量 & 100945.95 & \multicolumn{2}{c}{全天购电费（元）} & \multicolumn{2}{r}{77020.40}\\
\bottomrule
\end{tabularx}
\par\smallskip\parbox{\linewidth}{\footnotesize 注：问题2与4-2列示计划购电量；问题3与4-3列示最终生效的调整后常规购电量。全天购电量不含紧急购电，全天购电费含计划、调整与紧急费用。}
\end{table}

\begin{table}[htbp]\centering\small
\caption{问题4-2指定日期储能结果（沿用题目表2布局，单位：kWh）}
\label{tab:q4_2-battery}
\par\smallskip\textbf{2025-03-20}\par\smallskip
\begin{tabularx}{\linewidth}{lRRlRR}
\toprule
时间段 & 充电量 & 放电量 & 时间段 & 充电量 & 放电量\\
\midrule
00:00--04:00 & 3824.42 & 134.84 & 04:00--08:00 & 1669.83 & 5857.47\\
08:00--12:00 & 6146.04 & 2777.38 & 12:00--16:00 & 3907.13 & 420.71\\
16:00--20:00 & 174.78 & 6369.00 & 20:00--24:00 & 5798.62 & 2212.03\\
\midrule
0:00储电量 & \multicolumn{2}{r}{6834.43} & 24:00储电量 & \multicolumn{2}{r}{6457.15}\\
\bottomrule
\end{tabularx}
\par\smallskip\textbf{2025-06-21}\par\smallskip
\begin{tabularx}{\linewidth}{lRRlRR}
\toprule
时间段 & 充电量 & 放电量 & 时间段 & 充电量 & 放电量\\
\midrule
00:00--04:00 & 963.51 & 3780.66 & 04:00--08:00 & 1400.69 & 1719.95\\
08:00--12:00 & 6932.04 & 0.00 & 12:00--16:00 & 0.00 & 0.00\\
16:00--20:00 & 0.00 & 4853.45 & 20:00--24:00 & 7501.98 & 3091.92\\
\midrule
0:00储电量 & \multicolumn{2}{r}{8545.17} & 24:00储电量 & \multicolumn{2}{r}{8723.59}\\
\bottomrule
\end{tabularx}
\par\smallskip\textbf{2025-09-23}\par\smallskip
\begin{tabularx}{\linewidth}{lRRlRR}
\toprule
时间段 & 充电量 & 放电量 & 时间段 & 充电量 & 放电量\\
\midrule
00:00--04:00 & 5330.27 & 22.20 & 04:00--08:00 & 206.94 & 5820.07\\
08:00--12:00 & 5946.19 & 1896.62 & 12:00--16:00 & 3587.87 & 530.87\\
16:00--20:00 & 163.52 & 6562.31 & 20:00--24:00 & 5525.15 & 1842.66\\
\midrule
0:00储电量 & \multicolumn{2}{r}{6002.76} & 24:00储电量 & \multicolumn{2}{r}{6159.23}\\
\bottomrule
\end{tabularx}
\par\smallskip\textbf{2025-12-21}\par\smallskip
\begin{tabularx}{\linewidth}{lRRlRR}
\toprule
时间段 & 充电量 & 放电量 & 时间段 & 充电量 & 放电量\\
\midrule
00:00--04:00 & 4499.29 & 40.16 & 04:00--08:00 & 0.00 & 741.76\\
08:00--12:00 & 5940.70 & 6486.91 & 12:00--16:00 & 5724.46 & 2220.11\\
16:00--20:00 & 0.00 & 5500.96 & 20:00--24:00 & 5818.48 & 2842.53\\
\midrule
0:00储电量 & \multicolumn{2}{r}{6795.26} & 24:00储电量 & \multicolumn{2}{r}{6766.10}\\
\bottomrule
\end{tabularx}
\end{table}

\begin{table}[htbp]\centering\small
\caption{问题4-2指定日期紧急购电（题目表3布局，单位：kWh）}
\label{tab:q4_2-emergency}
\setlength{\tabcolsep}{3pt}
\begin{tabularx}{\linewidth}{lRlRlRlR}\toprule
\multicolumn{2}{c}{2025-03-20} & \multicolumn{2}{c}{2025-06-21} & \multicolumn{2}{c}{2025-09-23} & \multicolumn{2}{c}{2025-12-21}\\
时间段 & 购电量 & 时间段 & 购电量 & 时间段 & 购电量 & 时间段 & 购电量\\\midrule
无 & 0.00 & 无 & 0.00 & 18:20--18:40 & 3.91 & 无 & 0.00\\
 &   &  &   & 20:30--20:40 & 357.19 &  &  \\
 &   &  &   & 20:50--21:10 & 90.05 &  &  \\
 &   &  &   & 21:20--21:50 & 34.20 &  &  \\
\midrule
合计 & 0.00 & 合计 & 0.00 & 合计 & 485.35 & 合计 & 0.00\\
\bottomrule\end{tabularx}\end{table}
\begin{table}[htbp]\centering\small
\caption{问题4-2指定日期计划、调整与费用分解}
\begin{tabularx}{\linewidth}{lRRRRR}\toprule
日期 & 计划量/kWh & 生效量/kWh & 计划费/元 & 调整费/元 & 紧急费/元\\\midrule
2025-03-20 & 67869.67 & 67869.67 & 42207.16 & 0.00 & 0.00\\
2025-06-21 & 40675.72 & 40675.72 & 21012.94 & 0.00 & 0.00\\
2025-09-23 & 69762.68 & 69762.68 & 45194.04 & 0.00 & 3115.84\\
2025-12-21 & 100945.95 & 100945.95 & 77020.40 & 0.00 & 0.00\\
\bottomrule\end{tabularx}\end{table}
\FloatBarrier
\begin{table}[htbp]\centering\small
\caption{问题4-3指定日期购电结果（沿用题目表1布局，电量单位：kWh）}
\label{tab:q4_3-purchase}
\par\smallskip\textbf{2025-03-20}\par\smallskip
\begin{tabularx}{\linewidth}{lRlRlR}
\toprule
时间段 & 购电量 & 时间段 & 购电量 & 时间段 & 购电量\\
\midrule
10:00--10:10 & 0.00 & 12:00--12:10 & 624.29 & 14:00--14:10 & 0.00\\
16:00--16:10 & 671.04 & 18:00--18:10 & 704.20 & 20:00--20:10 & 743.27\\
\midrule
全天购电量 & 71574.04 & \multicolumn{2}{c}{全天购电费（元）} & \multicolumn{2}{r}{45107.31}\\
\bottomrule
\end{tabularx}
\par\smallskip\textbf{2025-06-21}\par\smallskip
\begin{tabularx}{\linewidth}{lRlRlR}
\toprule
时间段 & 购电量 & 时间段 & 购电量 & 时间段 & 购电量\\
\midrule
10:00--10:10 & 0.00 & 12:00--12:10 & 0.00 & 14:00--14:10 & 0.00\\
16:00--16:10 & 260.52 & 18:00--18:10 & 495.32 & 20:00--20:10 & 0.00\\
\midrule
全天购电量 & 37518.90 & \multicolumn{2}{c}{全天购电费（元）} & \multicolumn{2}{r}{19308.93}\\
\bottomrule
\end{tabularx}
\par\smallskip\textbf{2025-09-23}\par\smallskip
\begin{tabularx}{\linewidth}{lRlRlR}
\toprule
时间段 & 购电量 & 时间段 & 购电量 & 时间段 & 购电量\\
\midrule
10:00--10:10 & 22.75 & 12:00--12:10 & 461.07 & 14:00--14:10 & 54.35\\
16:00--16:10 & 611.01 & 18:00--18:10 & 0.00 & 20:00--20:10 & 822.11\\
\midrule
全天购电量 & 73431.36 & \multicolumn{2}{c}{全天购电费（元）} & \multicolumn{2}{r}{47890.72}\\
\bottomrule
\end{tabularx}
\par\smallskip\textbf{2025-12-21}\par\smallskip
\begin{tabularx}{\linewidth}{lRlRlR}
\toprule
时间段 & 购电量 & 时间段 & 购电量 & 时间段 & 购电量\\
\midrule
10:00--10:10 & 0.00 & 12:00--12:10 & 1054.19 & 14:00--14:10 & 0.00\\
16:00--16:10 & 877.48 & 18:00--18:10 & 744.23 & 20:00--20:10 & 0.00\\
\midrule
全天购电量 & 101918.89 & \multicolumn{2}{c}{全天购电费（元）} & \multicolumn{2}{r}{77624.69}\\
\bottomrule
\end{tabularx}
\par\smallskip\parbox{\linewidth}{\footnotesize 注：问题2与4-2列示计划购电量；问题3与4-3列示最终生效的调整后常规购电量。全天购电量不含紧急购电，全天购电费含计划、调整与紧急费用。}
\end{table}

\begin{table}[htbp]\centering\small
\caption{问题4-3指定日期储能结果（沿用题目表2布局，单位：kWh）}
\label{tab:q4_3-battery}
\par\smallskip\textbf{2025-03-20}\par\smallskip
\begin{tabularx}{\linewidth}{lRRlRR}
\toprule
时间段 & 充电量 & 放电量 & 时间段 & 充电量 & 放电量\\
\midrule
00:00--04:00 & 4119.21 & 134.84 & 04:00--08:00 & 825.23 & 5868.37\\
08:00--12:00 & 5765.88 & 2777.38 & 12:00--16:00 & 5187.95 & 254.72\\
16:00--20:00 & 28.20 & 6553.40 & 20:00--24:00 & 5792.33 & 1970.23\\
\midrule
0:00储电量 & \multicolumn{2}{r}{6530.78} & 24:00储电量 & \multicolumn{2}{r}{6567.77}\\
\bottomrule
\end{tabularx}
\par\smallskip\textbf{2025-06-21}\par\smallskip
\begin{tabularx}{\linewidth}{lRRlRR}
\toprule
时间段 & 充电量 & 放电量 & 时间段 & 充电量 & 放电量\\
\midrule
00:00--04:00 & 675.55 & 4135.87 & 04:00--08:00 & 1424.39 & 1719.95\\
08:00--12:00 & 8481.89 & 0.00 & 12:00--16:00 & 0.00 & 0.00\\
16:00--20:00 & 0.00 & 5210.22 & 20:00--24:00 & 6461.65 & 2485.90\\
\midrule
0:00储电量 & \multicolumn{2}{r}{7782.84} & 24:00储电量 & \multicolumn{2}{r}{8064.24}\\
\bottomrule
\end{tabularx}
\par\smallskip\textbf{2025-09-23}\par\smallskip
\begin{tabularx}{\linewidth}{lRRlRR}
\toprule
时间段 & 充电量 & 放电量 & 时间段 & 充电量 & 放电量\\
\midrule
00:00--04:00 & 5172.38 & 22.20 & 04:00--08:00 & 177.31 & 5680.26\\
08:00--12:00 & 6536.73 & 1893.23 & 12:00--16:00 & 3099.44 & 353.23\\
16:00--20:00 & 96.08 & 6443.21 & 20:00--24:00 & 5603.57 & 2259.51\\
\midrule
0:00储电量 & \multicolumn{2}{r}{6144.86} & 24:00储电量 & \multicolumn{2}{r}{6260.00}\\
\bottomrule
\end{tabularx}
\par\smallskip\textbf{2025-12-21}\par\smallskip
\begin{tabularx}{\linewidth}{lRRlRR}
\toprule
时间段 & 充电量 & 放电量 & 时间段 & 充电量 & 放电量\\
\midrule
00:00--04:00 & 4436.14 & 40.16 & 04:00--08:00 & 735.70 & 1384.74\\
08:00--12:00 & 6412.06 & 6391.52 & 12:00--16:00 & 5218.63 & 2240.52\\
16:00--20:00 & 0.00 & 5490.42 & 20:00--24:00 & 5769.36 & 2842.53\\
\midrule
0:00储电量 & \multicolumn{2}{r}{6852.09} & 24:00储电量 & \multicolumn{2}{r}{6733.60}\\
\bottomrule
\end{tabularx}
\end{table}

\begin{table}[htbp]\centering\small
\caption{问题4-3指定日期紧急购电（题目表3布局，单位：kWh）}
\label{tab:q4_3-emergency}
\setlength{\tabcolsep}{3pt}
\begin{tabularx}{\linewidth}{lRlRlRlR}\toprule
\multicolumn{2}{c}{2025-03-20} & \multicolumn{2}{c}{2025-06-21} & \multicolumn{2}{c}{2025-09-23} & \multicolumn{2}{c}{2025-12-21}\\
时间段 & 购电量 & 时间段 & 购电量 & 时间段 & 购电量 & 时间段 & 购电量\\\midrule
无 & 0.00 & 无 & 0.00 & 18:30--18:40 & 2.38 & 无 & 0.00\\
 &   &  &   & 20:10--20:20 & 10.23 &  &  \\
\midrule
合计 & 0.00 & 合计 & 0.00 & 合计 & 12.61 & 合计 & 0.00\\
\bottomrule\end{tabularx}\end{table}
\begin{table}[htbp]\centering\small
\caption{问题4-3指定日期计划、调整与费用分解}
\begin{tabularx}{\linewidth}{lRRRRR}\toprule
日期 & 计划量/kWh & 生效量/kWh & 计划费/元 & 调整费/元 & 紧急费/元\\\midrule
2025-03-20 & 71574.04 & 71574.04 & 45107.31 & 0.00 & 0.00\\
2025-06-21 & 37505.63 & 37518.90 & 19288.01 & 20.91 & 0.00\\
2025-09-23 & 73403.01 & 73431.36 & 47750.49 & 57.78 & 82.45\\
2025-12-21 & 101918.89 & 101918.89 & 77624.69 & 0.00 & 0.00\\
\bottomrule\end{tabularx}\end{table}
\FloatBarrier

\subsubsection{跨月份表现与模型适用边界}

图\ref{fig:v2-monthly}将四组最终配置的实际费用与紧急量按月汇总：问题二与第四问日前版保留母模型，问题三与第四问滚动版采用全时点筛选。月度波动同时包含季节负荷、光伏、价格及历史误差变化，不能仅凭某月曲线把差异归因于价格预测。联合场景通过同日配对保持数据中已存在的价荷关系，但28个历史日仍可能遗漏稀少高价缺口，也不能据此保证跨年份泛化。

\begin{figure}[htbp]\centering
\includegraphics[width=0.95\linewidth]{figures/results/v2_monthly.pdf}
\caption{问题二至四四组最终配置的月度实际费用与紧急购电量。日前两组采用母模型，滚动两组采用全时点筛选；每月数值由2--12月逐日执行账单和电量求和，非优化器辅助目标，不同方案的预测来源、筛选设置与初态差异需结合正文解释。}
\label{fig:v2-monthly}
\end{figure}

本问的机制结论是：未知价格使场景的经济权重随价格改变，故采购应联合考虑价荷而非只拟合需求均值；储能和滚动修订再处理跨时段与信息更新。配对情景下的同规则候选回放支持在本次开发期为滚动版采用全时点筛选，但日前版没有获益，紧急补购和尾部费用也可能增加。严格多阶段非预见性、未来修订机会的显式价值与完整跨日随机结算仍属于未实现的增强，不能将当前有限候选的费用最低结论推广为未来全局最优保证。
\FloatBarrier

\FloatBarrier
\section{模型评价与结论}

\subsection{四问结果与模型特点}

四个问题在相同物理状态方程与结算定义下递进。问题一的确定性LP解释储能如何在容量和效率约束下转移能量；问题二由五倍补购机制推导采购倾向，再用残差场景处理不确定性；问题三在合法交易时点按固定执行规则回放并筛选购电候选；问题四将同日价格残差与净负荷残差联合使用。最终配置为日前两组保留母模型、滚动两组在零时与日内均启用候选筛选，汇总如下。

\begin{table}[htbp]\centering\small
\caption{最终分问配置的334天开发期闭环评价结果}
\label{tab:annual_summary}
\begin{tabularx}{\linewidth}{LRRRR}\toprule
指标 & 问题2 & 问题3 & 问题4日前版 & 问题4滚动版\\\midrule
计划费用/万元 & 1350.80 & 1350.31 & 1428.01 & 1426.52\\
调整费用/万元 & 0.00 & 24.87 & 0.00 & 28.31\\
紧急费用/万元 & 50.48 & 11.97 & 51.40 & 12.19\\
实际总费用/万元 & 1401.28 & 1387.15 & 1479.41 & 1467.02\\
紧急电量/kWh & 85017.69 & 21029.90 & 80544.00 & 21293.41\\
紧急购电天数 & 113 & 119 & 106 & 116\\
紧急购电连续时段数 & 205 & 204 & 195 & 193\\
日总费用CVaR95/元 & 70731.74 & 65552.76 & 83044.31 & 77575.87\\
最高日费用/元 & 89306.82 & 67941.61 & 108080.37 & 81797.18\\
2月1日期初库存/kWh & 8875.19 & 8814.60 & 9251.54 & 9030.40\\
12月31日期末库存/kWh & 6240.32 & 6240.32 & 6269.86 & 6274.99\\
\bottomrule\end{tabularx}
\par\smallskip\parbox{\linewidth}{\footnotesize 问题2与问题4日前版沿用母模型；问题3与问题4滚动版在零时及日内均启用候选筛选。预测器与实时执行器沿用原冻结配置，筛选适用问题与时点依据开发期比较确定。实际费用不含候选评分的终端偏离项；CVaR95为实际日总账单的经验尾部均值。}
\end{table}


问题一费用为35126.95元，相对无储能节费26.90\%；问题二、三评价期主费用分别为1401.28万元、1387.15万元，动态电价日前版与滚动版分别为1479.41万元、1467.02万元。第三问与第四问滚动版的未筛选母模型费用分别为1397.58万元、1476.92万元，全时点筛选分别节省104339.41元和98994.66元，降幅为0.747\%和0.670\%。在同一问题四种配置末库存一致的条件下，全时点筛选取得本次比较中的最低账单。

这一收益伴随紧急补购增加：两滚动组分别由12647.76增至21029.90 kWh、由11950.68增至21293.41 kWh。计划费用的节省超过新增调整费用及紧急费用，并不表示所有运行指标同时改善。日前两组启用筛选反而分别增加0.161\%和0.428\%费用，因此保留母模型。信息来源、执行器和退款解释的原有对照均维持未筛选母模型配置，用于解释相应机制，不冒充最终增强算法的同条件消融。

主结果逐段满足供需平衡、储电上下界、功率限制、充放电互斥及日际连续性。未来实际值与未来发布预报的扰动没有改变已执行前缀。所有费用与时段分组由同一结果源生成，保留完整源代码、原模板及修订前版本，便于追溯图表中的每个数值。

\subsection{适用边界与可继续改进之处}

第一，只有一个年度且一月校准窗口较短；预测器与实时执行器沿用原冻结配置，筛选的适用问题和启用时点依据开发期结果确定，不能声称跨年份泛化或从未接触的独立测试。第二，预测与控制采用有限候选的坐标搜索，28场景、额外风险与备用关闭等设置保持不变；全时点筛选只在所比较的四种配置中费用最低，并未穷举全部可行策略。第三，候选回放沿用实际执行规则，部分改善了共享参考与实际控制的评估衔接，但回放固定当前合同至午夜，没有模拟后续尚未发布的预报及合同修订，不是完整多阶段随机最优控制。第四，终端偏离罚项仅用于优化和候选评分，不计入实际账单；模型未加入题目以外的网络限额、设备寿命退化和交易约束。

扩大候选集只能保证相同情景、初态和评分规则下的最小估计费用不增加，不能保证未来实际费用下降。本次保留具有消融依据的滚动采购增强，不再据开发期结果叠加风险、备用或新的实时控制器；其适用性仍需在其他未参与选型的时段或年度检验。


\FloatBarrier
\CumcmAIUsed{建模方案整理、预测与调度程序生成、数据计算、文献检索核验、图表生成及论文初稿撰写}
% 文献元数据核验日期：2026-09-10。
% DOI 注册元数据与出版社页面的核验记录见 ../文献与模型对应.md。
% 本文件直接使用 % 文献元数据核验日期：2026-09-10。
% DOI 注册元数据与出版社页面的核验记录见 ../文献与模型对应.md。
% 本文件直接使用 % 文献元数据核验日期：2026-09-10。
% DOI 注册元数据与出版社页面的核验记录见 ../文献与模型对应.md。
% 本文件直接使用 \input{sections/references}，无需 BibTeX。
\begin{thebibliography}{99}

\bibitem{kong2016}
Kong X, Bai L, Hu Q, et al.
Day-ahead optimal scheduling method for grid-connected microgrid based on energy storage control strategy[J].
Journal of Modern Power Systems and Clean Energy, 2016, 4(4): 648--658.
DOI: \url{https://doi.org/10.1007/s40565-016-0245-0}.

\bibitem{silva2020}
Silva V A, Aoki A R, Lambert-Torres G.
Optimal day-ahead scheduling of microgrids with battery energy storage system[J].
Energies, 2020, 13(19): 5188.
DOI: \url{https://doi.org/10.3390/en13195188}.

\bibitem{abunima2022}
Abunima H, Park W-H, Glick M B, et al.
Two-stage stochastic optimization for operating a renewable-based microgrid[J].
Applied Energy, 2022, 325: 119848.
DOI: \url{https://doi.org/10.1016/j.apenergy.2022.119848}.

\bibitem{marino2018}
Marino C, Quddus M A, Marufuzzaman M, et al.
A chance-constrained two-stage stochastic programming model for reliable microgrid operations under power demand uncertainty[J].
Sustainable Energy, Grids and Networks, 2018, 13: 66--77.
DOI: \url{https://doi.org/10.1016/j.segan.2017.12.007}.

\bibitem{grover2018}
Grover-Silva E, Heleno M, Mashayekh S, et al.
A stochastic optimal power flow for scheduling flexible resources in microgrids operation[J].
Applied Energy, 2018, 229: 201--208.
DOI: \url{https://doi.org/10.1016/j.apenergy.2018.07.114}.

\bibitem{palma2013}
Palma-Behnke R, Benavides C, Lanas F, et al.
A microgrid energy management system based on the rolling horizon strategy[J].
IEEE Transactions on Smart Grid, 2013, 4(2): 996--1006.
DOI: \url{https://doi.org/10.1109/TSG.2012.2231440}.

\bibitem{silvente2015}
Silvente J, Kopanos G M, Pistikopoulos E N, et al.
A rolling horizon optimization framework for the simultaneous energy supply and demand planning in microgrids[J].
Applied Energy, 2015, 155: 485--501.
DOI: \url{https://doi.org/10.1016/j.apenergy.2015.05.090}.

\bibitem{li2016}
Li Z, Zang C, Zeng P, et al.
Combined two-stage stochastic programming and receding horizon control strategy for microgrid energy management considering uncertainty[J].
Energies, 2016, 9(7): 499.
DOI: \url{https://doi.org/10.3390/en9070499}.

\bibitem{jani2022}
Jani A, Karimi H, Jadid S.
Two-layer stochastic day-ahead and real-time energy management of networked microgrids considering integration of renewable energy resources[J].
Applied Energy, 2022, 323: 119630.
DOI: \url{https://doi.org/10.1016/j.apenergy.2022.119630}.

\bibitem{li2020}
Li J, Tan C, Ren Z, et al.
A two-stage optimal scheduling model of microgrid based on chance-constrained programming in spot markets[J].
Processes, 2020, 8(1): 107.
DOI: \url{https://doi.org/10.3390/pr8010107}.

\bibitem{rockafellar2000}
Rockafellar R T, Uryasev S.
Optimization of conditional value-at-risk[J].
The Journal of Risk, 2000, 2(3): 21--41.
DOI: \url{https://doi.org/10.21314/JOR.2000.038}.

\bibitem{elmachtoub2017spo}
Elmachtoub A N, Grigas P.
Smart ``Predict, then Optimize''[EB/OL].
arXiv:1710.08005, 2017 (v5 updated 2020-11-19) [2026-09-10].
\url{https://arxiv.org/abs/1710.08005}.
DOI: \url{https://doi.org/10.48550/arXiv.1710.08005}.

\end{thebibliography}
，无需 BibTeX。
\begin{thebibliography}{99}

\bibitem{kong2016}
Kong X, Bai L, Hu Q, et al.
Day-ahead optimal scheduling method for grid-connected microgrid based on energy storage control strategy[J].
Journal of Modern Power Systems and Clean Energy, 2016, 4(4): 648--658.
DOI: \url{https://doi.org/10.1007/s40565-016-0245-0}.

\bibitem{silva2020}
Silva V A, Aoki A R, Lambert-Torres G.
Optimal day-ahead scheduling of microgrids with battery energy storage system[J].
Energies, 2020, 13(19): 5188.
DOI: \url{https://doi.org/10.3390/en13195188}.

\bibitem{abunima2022}
Abunima H, Park W-H, Glick M B, et al.
Two-stage stochastic optimization for operating a renewable-based microgrid[J].
Applied Energy, 2022, 325: 119848.
DOI: \url{https://doi.org/10.1016/j.apenergy.2022.119848}.

\bibitem{marino2018}
Marino C, Quddus M A, Marufuzzaman M, et al.
A chance-constrained two-stage stochastic programming model for reliable microgrid operations under power demand uncertainty[J].
Sustainable Energy, Grids and Networks, 2018, 13: 66--77.
DOI: \url{https://doi.org/10.1016/j.segan.2017.12.007}.

\bibitem{grover2018}
Grover-Silva E, Heleno M, Mashayekh S, et al.
A stochastic optimal power flow for scheduling flexible resources in microgrids operation[J].
Applied Energy, 2018, 229: 201--208.
DOI: \url{https://doi.org/10.1016/j.apenergy.2018.07.114}.

\bibitem{palma2013}
Palma-Behnke R, Benavides C, Lanas F, et al.
A microgrid energy management system based on the rolling horizon strategy[J].
IEEE Transactions on Smart Grid, 2013, 4(2): 996--1006.
DOI: \url{https://doi.org/10.1109/TSG.2012.2231440}.

\bibitem{silvente2015}
Silvente J, Kopanos G M, Pistikopoulos E N, et al.
A rolling horizon optimization framework for the simultaneous energy supply and demand planning in microgrids[J].
Applied Energy, 2015, 155: 485--501.
DOI: \url{https://doi.org/10.1016/j.apenergy.2015.05.090}.

\bibitem{li2016}
Li Z, Zang C, Zeng P, et al.
Combined two-stage stochastic programming and receding horizon control strategy for microgrid energy management considering uncertainty[J].
Energies, 2016, 9(7): 499.
DOI: \url{https://doi.org/10.3390/en9070499}.

\bibitem{jani2022}
Jani A, Karimi H, Jadid S.
Two-layer stochastic day-ahead and real-time energy management of networked microgrids considering integration of renewable energy resources[J].
Applied Energy, 2022, 323: 119630.
DOI: \url{https://doi.org/10.1016/j.apenergy.2022.119630}.

\bibitem{li2020}
Li J, Tan C, Ren Z, et al.
A two-stage optimal scheduling model of microgrid based on chance-constrained programming in spot markets[J].
Processes, 2020, 8(1): 107.
DOI: \url{https://doi.org/10.3390/pr8010107}.

\bibitem{rockafellar2000}
Rockafellar R T, Uryasev S.
Optimization of conditional value-at-risk[J].
The Journal of Risk, 2000, 2(3): 21--41.
DOI: \url{https://doi.org/10.21314/JOR.2000.038}.

\bibitem{elmachtoub2017spo}
Elmachtoub A N, Grigas P.
Smart ``Predict, then Optimize''[EB/OL].
arXiv:1710.08005, 2017 (v5 updated 2020-11-19) [2026-09-10].
\url{https://arxiv.org/abs/1710.08005}.
DOI: \url{https://doi.org/10.48550/arXiv.1710.08005}.

\end{thebibliography}
，无需 BibTeX。
\begin{thebibliography}{99}

\bibitem{kong2016}
Kong X, Bai L, Hu Q, et al.
Day-ahead optimal scheduling method for grid-connected microgrid based on energy storage control strategy[J].
Journal of Modern Power Systems and Clean Energy, 2016, 4(4): 648--658.
DOI: \url{https://doi.org/10.1007/s40565-016-0245-0}.

\bibitem{silva2020}
Silva V A, Aoki A R, Lambert-Torres G.
Optimal day-ahead scheduling of microgrids with battery energy storage system[J].
Energies, 2020, 13(19): 5188.
DOI: \url{https://doi.org/10.3390/en13195188}.

\bibitem{abunima2022}
Abunima H, Park W-H, Glick M B, et al.
Two-stage stochastic optimization for operating a renewable-based microgrid[J].
Applied Energy, 2022, 325: 119848.
DOI: \url{https://doi.org/10.1016/j.apenergy.2022.119848}.

\bibitem{marino2018}
Marino C, Quddus M A, Marufuzzaman M, et al.
A chance-constrained two-stage stochastic programming model for reliable microgrid operations under power demand uncertainty[J].
Sustainable Energy, Grids and Networks, 2018, 13: 66--77.
DOI: \url{https://doi.org/10.1016/j.segan.2017.12.007}.

\bibitem{grover2018}
Grover-Silva E, Heleno M, Mashayekh S, et al.
A stochastic optimal power flow for scheduling flexible resources in microgrids operation[J].
Applied Energy, 2018, 229: 201--208.
DOI: \url{https://doi.org/10.1016/j.apenergy.2018.07.114}.

\bibitem{palma2013}
Palma-Behnke R, Benavides C, Lanas F, et al.
A microgrid energy management system based on the rolling horizon strategy[J].
IEEE Transactions on Smart Grid, 2013, 4(2): 996--1006.
DOI: \url{https://doi.org/10.1109/TSG.2012.2231440}.

\bibitem{silvente2015}
Silvente J, Kopanos G M, Pistikopoulos E N, et al.
A rolling horizon optimization framework for the simultaneous energy supply and demand planning in microgrids[J].
Applied Energy, 2015, 155: 485--501.
DOI: \url{https://doi.org/10.1016/j.apenergy.2015.05.090}.

\bibitem{li2016}
Li Z, Zang C, Zeng P, et al.
Combined two-stage stochastic programming and receding horizon control strategy for microgrid energy management considering uncertainty[J].
Energies, 2016, 9(7): 499.
DOI: \url{https://doi.org/10.3390/en9070499}.

\bibitem{jani2022}
Jani A, Karimi H, Jadid S.
Two-layer stochastic day-ahead and real-time energy management of networked microgrids considering integration of renewable energy resources[J].
Applied Energy, 2022, 323: 119630.
DOI: \url{https://doi.org/10.1016/j.apenergy.2022.119630}.

\bibitem{li2020}
Li J, Tan C, Ren Z, et al.
A two-stage optimal scheduling model of microgrid based on chance-constrained programming in spot markets[J].
Processes, 2020, 8(1): 107.
DOI: \url{https://doi.org/10.3390/pr8010107}.

\bibitem{rockafellar2000}
Rockafellar R T, Uryasev S.
Optimization of conditional value-at-risk[J].
The Journal of Risk, 2000, 2(3): 21--41.
DOI: \url{https://doi.org/10.21314/JOR.2000.038}.

\bibitem{elmachtoub2017spo}
Elmachtoub A N, Grigas P.
Smart ``Predict, then Optimize''[EB/OL].
arXiv:1710.08005, 2017 (v5 updated 2020-11-19) [2026-09-10].
\url{https://arxiv.org/abs/1710.08005}.
DOI: \url{https://doi.org/10.48550/arXiv.1710.08005}.

\end{thebibliography}


\clearpage
\begin{appendices}
\section{支撑材料与结果生成链}

论文入口为\texttt{main.tex}，正文四至七章分别为四个问题，每章依次包含分析、模型建立和求解。原\texttt{example.tex}、类文件及样式文件保留。计算代码、原始模板和第一版源代码分别保存在\path{code/}、\path{originals/}、\path{versions/v1_before_settlement_revision/}；分问重排前的正文与代码另存\path{versions/v2_before_question_restructure/}。在项目目录使用\texttt{latexmk -xelatex main.tex}可编译当前论文。

\begin{table}[htbp]\centering\small
\caption{修订稿计算入口与结果文件}\label{tab:files}
\begin{tabularx}{\linewidth}{p{0.43\linewidth}X}\toprule
文件 & 用途\\\midrule
\path{code/data_forecast.py} & 原始附件读取及因果预测基础\\
\path{code/optimize_dispatch.py} & 保留原版LP、残差场景与规则执行\\
\path{code/dispatch_v2.py} & 修订控制器、跨日续接与逐段结算\\
\path{code/calibrate_v2.py} & 一月费用选型、匹配参数与冻结配置\\
\path{code/run_main_v2.py} & 四组未筛选母模型结果\\
\path{code/outer_candidate_selector.py} & 候选生成与固定执行规则回放评分\\
\path{code/run_outer_candidate_experiment.py} & 母模型与全时点筛选对照\\
\path{code/run_outer_timing_ablation.py} & 四次筛选时点消融与配置冻结\\
\path{code/publish_final_strategy.py} & 最终配置汇总、主结果图表及附件数据\\
\path{code/experiments_v2.py} & 场景、风险、备用、反馈及控制对照\\
\path{code/diagnose_execution_gap.py} & 原版全决策点规划与规则回放诊断\\
\path{code/diagnose_inner_v2.py} & 指定四日的同合同两控制器诊断\\
\path{code/time_index.py} & 物理时间、模板列序及分组映射\\
\path{code/make_paper_assets_v2.py} & 母模型数表、插图及原对照\\
\path{code/make_question_tables.py} & 分问编排中的预测与场景对照表\\
\path{code/make_visual_atlas.py} & 热力矩阵、瀑布图与场景阴影带\\
\path{code/scenario_convergence.py} & 真实场景缩减迭代记录与收敛分位带\\
\path{code/check_dispatch_v2.py} & 数值守恒、未来扰动及价格机制检查\\
\path{code/export_final_workbooks.mjs} & 从最终导出数据重填两本滚动工作簿\\
\path{code/verify_final_workbooks.py} & 五本最终附件的逐值、标签及汇总检查\\
\path{code/check_final_strategy_delivery.py} & 最终执行、图表、正文与附件独立核对\\
\path{code/run_revision.ps1} & 默认复用冻结配置的修订复算入口\\
\bottomrule\end{tabularx}
\end{table}

\texttt{selection\_v2.json}保存每个一月候选、费用、末库存、母模型配置及固定参数匹配值。\texttt{prepared\_v2\_*.npz}为已选预测输入，\texttt{*\_v2.npz}为334天母模型执行结果；旧\texttt{*\_full.npz}是原版365天记录，提供冻结的一月运行前缀。最终主结果中，问题二及第四问日前版仍读取母模型数组，问题三及第四问滚动版分别读取\path{outputs/outer_candidate_test/}中的\texttt{q3\_selected.npz}及\texttt{q4\_3\_selected.npz}。

筛选时点的四次回测、逐决策评分、配置及原始检查保存在\path{outputs/outer_timing_ablation/}；其中\texttt{final\_strategy.json}明确每问所选配置和执行数组路径。原运行清单保留当时输入的指纹，论文与附件随后获准更新，不回写旧清单伪装为原实验输入。当前发布文件及其来源另由\path{outputs/final_strategy/}中的清单记录。

每个主结果中，\texttt{plan}为零时合同，\texttt{adjusted}为每段最终有效常规量，\texttt{charge}、\texttt{discharge}为电池外侧实际电量，\texttt{energy}每一天含145个边界状态。母模型另以\texttt{reference\_energy}保存外层参考路径，最终两组滚动数组不包含该字段，故不报告混合含义的参考SOC偏差。\texttt{emergency}与\texttt{spill}分别为紧急与弃余电量；\texttt{plan\_cost}、\texttt{adjustment\_cost}及\texttt{emergency\_cost}之和为\texttt{total\_cost}。所有量保持完整精度，表内显示再按相应小数位舍入。

最终主执行数组汇入\path{outputs/final_strategy/results_export.json}，同目录的\texttt{summary.json}及\texttt{daily\_summary.csv}保存主指标、逐日费用和库存。工作簿与正文数表使用对应最终数组和统一索引，不从PDF抄录汇总数值。\path{data/results_export.json}保留母模型导出，不能用于覆盖最终滚动附件。原\path{experiments_v2.json/npz}及规划执行诊断仍保存母模型对照，不与新主案费用拼成未经运行的实验。

\section{时间映射与附件5填写}

附件5购电页保持原模板首标签“0:10--0:20”和末标签次日“0:00--0:10”。首末标签的原始文本在不同工作簿中略有差别，均逐字保留。以下列出源样本、物理模型与模板显示的对应，不据标签再平移数组：

\begin{center}\small
\begin{tabular}{llll}\toprule
数组索引 & 源样本右端点 & 物理时段 & 保留模板显示\\\midrule
$t=0$ & 00:10 & 00:00--00:10 & 0:10--0:20\\
$t=60$ & 10:10 & 10:00--10:10 & 10:10--10:20\\
$t=143$ & 次日00:00 & 23:50--24:00 & 次日0:00--0:10\\
\bottomrule\end{tabular}
\end{center}

具体原标签、Excel行列与四小时块均冻结于\path{data/time_mapping.json/csv}。问题一为纵向144行，其他购电页为横向144列；物理10:00--10:10从同一$t=60$取值。四小时储能区间依次汇总24段，紧急页将同一天相邻正紧急量合并，跨日不合并，无紧急量日期写“无”和0。源样本按右端点平均功率解释是明示建模假设，不将模板显示差异说成题目已给定的物理事实。

计划页的全天购电量及费用分别为$\sum_tg_t$及$\sum_tp_tg_t$；调整页填$a_t$而非$a_t-g_t$，全天费用填式\eqref{eq:settlement}的全部实际账单。没有调整页的问题二与问题四静态案，其完整账单由逐日CSV和正文给出。五本已填文件保存在\texttt{附件/附件5/}，空白原模板仅保存在\texttt{originals/附件5/}，交付时必须区别二者。

原母模型填报曾逐格核对288720个计划及调整数值，记录保留在\path{outputs/workbook_previews/}。本次按最终配置重写\texttt{result3.xlsx}及\texttt{result4-3.xlsx}，另外三本保持原方案；两本滚动附件的计划及调整区域共192384个数值，并重新核对储能聚合、紧急区间、日期与模板标签。当前文件指纹、导出与独立核对记录保存在\path{outputs/final_strategy/}，覆盖前文件另行备份。文件一致性检查不替代独立复现模型和人工论文审查。

\section{数值可行性与因果信息检查}

四个最终主案检查供给量不小于净需求；允许弃余时，使用扣除弃置量后的等式残差核对账本，不将优化约束改为强制全部电量使用。另检查SOC递推、日际连续性、1200--10800 kWh库存界、每段$5000/6$ kWh充放电上限及逐段互斥。计划费用始终按全部零时合同量结算，富余或未使用部分不抵扣。

母模型的未来扰动检查见\path{data/numerical_checks_v2.json}。新增候选模块还检查了逐场景回放与标量执行器一致性、日内保留当前合同、原计划结算基准、配对场景价格、终端偏离期望以及未观测实值和未发布预报扰动后已执行前缀不变；时点消融检查关闭筛选后仍求原场景LP。详细结果在两个外层实验目录的检查记录中。有限因果检查不等同于任意输入的形式证明，也不能使开发期选型变成独立测试。

以下节选显示外层交易与内层动作的职责。完整数值实现及参数保存在源文件中。

\begin{lstlisting}[language=Python,basicstyle=\ttfamily\footnotesize]
for t in range(144):
    if t in allowed_updates:
        if screening_enabled(question, t):
            candidates = candidate_plans(legal_forecast, observed_energy,
                                         original_day_ahead_plan, contract)
            contract = select_by_fixed_rule_replay(candidates,
                       same_historical_scenarios, observed_energy)
        else:
            contract = outer_plan(legal_forecast, observed_energy,
                                  original_day_ahead_plan)
    # Contract remains fixed inside the storage controller.
    action = controller(contract[t:], current_net, current_price,
                        latest_legal_future_forecast, observed_energy)
    execute_first_action(action)

plan_cost = sum(price * original_plan)
adjustment_cost = sum(price * (1.5 * maximum(final_plan-original_plan, 0)
                              + 0.5 * maximum(original_plan-final_plan, 0)))
emergency_cost = sum(5 * price * emergency)
total_cost = plan_cost + adjustment_cost + emergency_cost
\end{lstlisting}

\section{指定日期费用分解}

以下附表将正文指定日期的常规量及三项费用并列，便于与工作簿计划页、调整页及紧急页交叉定位。
\begin{table}[htbp]\centering\small
\caption{问题2指定日期计划、调整与费用分解}
\begin{tabularx}{\linewidth}{lRRRRR}\toprule
日期 & 计划量/kWh & 生效量/kWh & 计划费/元 & 调整费/元 & 紧急费/元\\\midrule
2025-03-20 & 67796.64 & 67796.64 & 41219.41 & 0.00 & 0.00\\
2025-06-21 & 40036.30 & 40036.30 & 23986.38 & 0.00 & 0.00\\
2025-09-23 & 69576.50 & 69576.50 & 43098.28 & 0.00 & 3285.98\\
2025-12-21 & 101008.53 & 101008.53 & 65454.44 & 0.00 & 0.00\\
\bottomrule\end{tabularx}\end{table}

\end{appendices}
\end{document}
